% SOURCE main.tex
\documentclass[11pt]{article}
\usepackage[T1]{fontenc}
\usepackage[utf8]{inputenc}
\usepackage{lmodern}
\usepackage[margin=1in]{geometry}
\usepackage{amsmath,amssymb,amsthm,mathtools,booktabs,array,microtype}
\usepackage{xurl}
\usepackage[colorlinks=true,linkcolor=blue,citecolor=blue,urlcolor=blue]{hyperref}
\hypersetup{pdftitle={Explicit Strong-Coupling Bounds and Vacuum-Adapted Coercivity for SU(N) Lattice Gauge Hamiltonians},pdfauthor={Carl P. Zimmerman}}
\newtheorem{theorem}{Theorem}[section]
\newtheorem{lemma}[theorem]{Lemma}
\newtheorem{proposition}[theorem]{Proposition}
\newtheorem{corollary}[theorem]{Corollary}
\theoremstyle{definition}
\newtheorem{definition}[theorem]{Definition}
\theoremstyle{remark}
\newtheorem{remark}[theorem]{Remark}
\DeclareMathOperator{\gap}{gap}
\DeclareMathOperator{\SU}{SU}
\DeclareMathOperator{\Tr}{Tr}
\DeclareMathOperator{\Var}{Var}
\DeclareMathOperator{\Ran}{Ran}
\DeclareMathOperator{\Dom}{Dom}
\newcommand{\dd}{\,\mathrm d}
\numberwithin{equation}{section}
\setlength{\emergencystretch}{2em}
\title{Explicit Strong-Coupling Bounds and\\
Vacuum-Adapted Coercivity for\\
$\SU(N)$ Lattice Gauge Hamiltonians}
\author{Carl P. Zimmerman\\[2pt]\small Briar Creek Tech}
\date{October 28, 2026 \quad Version 1.0}
\begin{document}
\maketitle
\begin{abstract}
We give quantitative spectral estimates for finite-lattice $\SU(N)$
gauge Hamiltonians and identify the additional estimates needed for a
continuum mass-gap argument. A sector-valued Kirkwood--Thomas expansion
extends a weak-interaction argument to bounded interactions on at most
$s$ sites, with constants independent of the local Hilbert-space dimension.
For a specified normalization of the Kogut--Susskind Hamiltonian in three
spatial dimensions, this yields a finite-volume gap
$\Delta_{\mathrm{lat}}\ge 3g^2/16$ whenever $g^2\ge59.9$, uniformly in
$N\ge2$ and lattice size. We also give a counterexample to an intermediate
counting assertion in the original Bravyi--DiVincenzo--Loss preprint and
show that its conservative numerical constant survives the correction.
For a fixed block in the actual interacting finite-lattice vacuum, a
Feynman--Kac comparison yields a positive discarded-sector bound uniform
in the exterior volume. Finally, an exact Schur-complement comparison
adds inverse-gap costs when the induced coarse inner product is retained.
The first two spectral estimates are unconditional within their stated
finite-lattice hypotheses; their iteration into a nonlinear multiscale
construction is not established. We do not prove continuum Yang--Mills
existence, a continuum mass gap, or confinement. The accompanying exact
arithmetic checks certify specified constants and counterexamples, not
the full analytic arguments.
\end{abstract}

\section{Introduction and scope}\label{sec:introduction}

The Yang--Mills existence and mass-gap problem asks for a nontrivial
four-dimensional quantum field theory for every compact simple gauge
group, together with a strictly positive physical mass gap and the
required field-theoretic axioms~\cite{JW}. A positive eigenvalue gap for
one finite regulator is a much weaker assertion. Even a lattice-size
uniform gap at fixed strong coupling does not by itself control the
weak bare-coupling trajectory used in a continuum limit.

This paper isolates several finite-regulator results and an exact
conditional comparison. The principal quantitative result is a
dimension-independent extension of the creation-operator estimates of
Bravyi, DiVincenzo and Loss~\cite{BDL}. The extension is applied to a
Hamiltonian lattice gauge model of Kogut--Susskind type~\cite{KS}.
Weak-interaction stability of gapped product systems, including systems
with infinite-dimensional local spaces, is already established in
Yarotsky's work~\cite{Yarotsky}; the existence of a sufficiently
strong-coupling lattice gap is not claimed as a new phenomenon here.
Our emphasis is the explicit sector estimate, the resulting sufficient
constants, and careful separation of finite-lattice and continuum claims.
We make no claim that these constants are optimal or that their precise
form has not appeared elsewhere.

The second result concerns the actual ground-state measure, rather than
a product reference vacuum. A local Feynman--Kac comparison controls the
oscillation of the ground function over a fixed block independently of
the size of its exterior. Product-Haar Poincar\'e comparison then gives
coercivity on the sector discarded by conditional expectation. The
heat-kernel ingredient is the compact nonnegative-Ricci case of the
Li--Yau inequality~\cite{LiYau}. This result holds at every positive
coupling for a fixed finite block, but its constants do not remain
uniform when that block becomes an extensive elimination layer.

The final comparison records which operator must be controlled after
exact elimination. A Schur form alone omits the norm of its graph lift.
Keeping that norm gives a harmonic gap bound and additive inverse-gap
costs. We provide counterexamples to two tempting replacements of this
argument: ignoring the graph norm, and identifying quantum Schur
elimination with a transformation that preserves quadratic oscillator
Hamiltonians.

\subsection{What is proved, and what is conditional}

\begin{center}
\small
\begin{tabular}{@{}p{.25\textwidth}p{.32\textwidth}p{.35\textwidth}@{}}
\toprule
Result & Proved scope & Limitation\\
\midrule
Explicit strong-coupling stability &
Finite systems; constants uniform in size and on-site dimension &
No explicit-threshold infinite-volume construction or zero-spacing limit\\[4pt]
Gauge-Hamiltonian corollary &
Specified $\SU(N)$ normalization; all $N\ge2$ at fixed $g^2$ &
Uniformity is not at fixed 't Hooft coupling\\[4pt]
Actual-vacuum block coercivity &
Fixed finite block; all exterior volumes and configurations &
No uniform whole-layer or iterated coarse-operator estimate\\[4pt]
Schur gap transfer &
Exact operators and stated form domains, with their induced metrics &
Application to successive Yang--Mills scales remains conditional\\
\bottomrule
\end{tabular}
\end{center}

\subsection{Conventions and provenance}

Let $E$ be a finite set of oriented lattice links, with one group variable
per link, and let $\mathcal P$ be its plaquettes. Haar measure is normalized.
The kinematic Hilbert space and dimensionless Hamiltonian are
\begin{align}
 \mathcal H_E&=L^2(\SU(N)^E,\dd U), \qquad x=g^2>0,\quad N\ge2,\\
 H_E(x)&=\frac{x}{2}\sum_{\ell\in E}C_\ell+
          \frac{b_N}{x}\sum_{p\in\mathcal P}(1-T_p),&
 T_p&=\frac1N\operatorname{Re}\Tr U_p .\label{eq:hamiltonian}
\end{align}
The nonnegative Laplace--Beltrami operator $C_\ell=-\Delta_\ell$ uses the
Casimir normalization
\begin{equation}
 C_F=\frac{N^2-1}{2N},\qquad
 \inf\bigl(\operatorname{spec}C_\ell\setminus\{0\}\bigr)=C_F.
\end{equation}
The magnetic coefficient is an explicitly declared parameter.
The uniform strong-coupling statement assumes $0\le b_N\le2N$;
the finite-block theorem allows any fixed $b=b_N\ge0$. Changing conventions
requires transforming this parameter and the energy normalization.
If $a>0$ is the lattice spacing, physical excitation energies in units
$\hbar=c=1$ are eigenvalues of $(H_E-E_0)/a$.
The physical gauge-invariant space is the invariant subspace under the
local $\SU(N)$ gauge transformations. Both the kinematic space and this
restriction are specified whenever a gap is asserted.

The arguments assembled here were developed in September 2026 research
notes. The earlier spectral paper~\cite{Spectral} contains the
fixed-spacing qualitative strong-coupling application. The later
quantitative and vacuum-adapted notes are included in the October 5,
2026 versioned repository archive~\cite{Archive}. The present paper gives
a standalone exposition, exact supporting certificates, and an explicit
account of what is still missing. These dates document provenance;
they are not a claim of priority over other research.

\section{An explicit finite-volume strong-coupling bound}
\label{sc:section}

This section gives a quantitative version of stability of a gapped product
Hamiltonian under bounded local perturbations. The argument follows the
Kirkwood--Thomas method used by Bravyi, DiVincenzo, and Loss
\cite{BDL}, with Hilbert-space norms on excitation sectors in place of
absolute values of qubit coefficients. The existence of a sufficiently
small perturbative window, including for infinite-dimensional site spaces,
is established in the literature \cite{Yarotsky}. The purpose here is to
track constants for the specified Hamiltonian lattice gauge model.
Neither the method nor the existence of strong-coupling lattice gaps is
claimed to be new. The derivation and associated records originated in
the research archive \cite{Archive}; the earlier spectral-bounds paper
\cite{Spectral} gives the corresponding non-explicit stability application.

\subsection{A finite-system stability criterion}

Let $\mathcal L$ be a nonempty finite set of sites. At each site $u$, let
$\mathcal H_u$ be a separable Hilbert space and let $h_u$ be a nonnegative
self-adjoint operator with compact resolvent. Assume that $h_u$ has a
normalized simple zero eigenvector $\Omega_u$ and that, for a common
$\Delta>0$,
\begin{equation}
 h_u\big|_{\Omega_u^\perp}\ge\Delta.
 \label{sc:onsite-gap}
\end{equation}
The inequality is vacuous if $\mathcal H_u$ is one-dimensional. On
$\mathcal H=\bigotimes_{u\in\mathcal L}\mathcal H_u$, set
\begin{equation}
 H_0=\sum_{u\in\mathcal L}h_u,\qquad
 V=\sum_{p\in\mathcal P}V_p,\qquad
 H(\varepsilon)=H_0+\varepsilon V.
 \label{sc:general-H}
\end{equation}
Here $\mathcal P$ is a finite indexed family of nonempty site sets, each
of cardinality at most an integer $s\ge1$. Each $V_p$ is self-adjoint, supported on
$p$, and has norm at most $J$. Every site belongs to at most $D$
members of $\mathcal P$, counted with their multiplicities. We may take
$D,J>0$; when there is no perturbation the conclusion below is immediate.
The coupling $\varepsilon$ is real.

Define the dimensionless parameter and positive constants
\begin{align}
 \lambda&=\frac{|\varepsilon|DJ}{\Delta},
 \label{sc:lambda}\\
 \gamma_0(s)^2&=\sum_{t=1}^{s}\frac{\binom{s-1}{t-1}}{t^2},
 &\gamma_1(s)^2&=(s+1)^2\sum_{t=0}^{s}
                        \frac{\binom{s}{t}}{(t+1)^2},
 \label{sc:gammas}\\
 \beta_0(s)&=\gamma_0(s),
 &\beta_k(s)&=2^k\bigl(\gamma_0(s)s^k+
                         k\gamma_1(s)s^{k-1}\bigr)
                         \quad(1\le k\le2s),
 \label{sc:betas}\\
 q_s(r)&=\sum_{k=0}^{2s}\frac{\beta_k(s)}{k!}r^k.
 \label{sc:majorant}
\end{align}

\begin{theorem}[Dimension-independent finite-volume stability]
\label{sc:stability}
Suppose there is a number $r>0$ such that
\begin{equation}
 \lambda q_s(r)\le r,
 \qquad 2\lambda q_s'(r)<1.
 \label{sc:criterion}
\end{equation}
Then $H(\varepsilon)$ has a simple ground eigenvalue $E_0$ and
\begin{equation}
 \operatorname{spec}H(\varepsilon)
 \subset\{E_0\}\cup[E_0+\Delta/2,\infty).
 \label{sc:gap}
\end{equation}
The criterion depends on $s,D,J,\Delta$ and $\varepsilon$, but not on
the number of sites, their Hilbert-space dimensions, or the upper
onsite spectrum.
\end{theorem}

We first prove the theorem when all site spaces are finite-dimensional.
The compact-resolvent extension is proved at the end of this subsection.

For $M\subset\mathcal L$, let
$\mathcal K_M=\bigotimes_{u\in M}\Omega_u^\perp$ and
$\Omega_M=\bigotimes_{u\in M}\Omega_u$. The full Hilbert space is the
orthogonal sum of the sectors
$\mathcal K_M\otimes\Omega_{\mathcal L\setminus M}$.
We identify a sector with $\mathcal K_M$, and write $\Pi_M$ for the
corresponding projection followed by this identification. In particular,
\begin{equation}
 \sum_{M\subset\mathcal L}\|\Pi_M\xi\|^2=\|\xi\|^2.
 \label{sc:orthogonal-sectors}
\end{equation}
Let $\Omega=\bigotimes_{u\in\mathcal L}\Omega_u$.
For $c_M\in\mathcal K_M$, $M\ne\varnothing$, define
\begin{equation}
 A[c_M]=|c_M\rangle\langle\Omega_M|\otimes
                  I_{\mathcal L\setminus M},
 \qquad
 C=\sum_{M\ne\varnothing}A[c_M],
 \qquad
 \|C\|_1=\max_{u\in\mathcal L}
                 \sum_{M\ni u}\|c_M\|.
 \label{sc:creation-norm}
\end{equation}
The decomposition of $C$ is unique, as is seen by applying it to $\Omega$.
We call $C$ a creation operator. The coefficients $c_M$ are vectors,
including when the excitations within a sector are entangled.

Two such elementary operators commute. If their supports intersect,
their product in either order is zero: a vacuum bra on an overlapping
site annihilates an excited factor. If their supports are disjoint,
they commute by tensor factorization. Consequently all creation operators
commute, and every product of more than $|\mathcal L|$ elementary creation
operators vanishes. Their exponentials are therefore finite polynomials.
Also, $\|A[c_M]\|=\|c_M\|$.

Put $H_0^M=\sum_{u\in M}h_u$ on $\mathcal K_M$. Then
\begin{equation}
 H_0^M\ge\Delta|M|,
 \qquad [H_0,A[c_M]]=A[H_0^M c_M].
 \label{sc:free-commutator}
\end{equation}
Thus a second commutator of $H_0$ with creation operators is zero.
We use $\operatorname{ad}_C(X)=[C,X]$.

\begin{lemma}[Support and commutator depth]
\label{sc:support}
Fix $p\in\mathcal P$ and elementary creation operators
$A[c_{M_1}],\ldots,A[c_{M_k}]$. Set $N_*=\bigcup_{i=1}^kM_i$ and
\[
 X=\operatorname{ad}_{A[c_{M_1}]}\cdots
                   \operatorname{ad}_{A[c_{M_k}]}(V_p).
\]
If $X\ne0$, each $M_i$ intersects $p$. If $\Pi_MX\Omega\ne0$, then
\begin{equation}
 N_*\setminus p\subset M\subset N_*\cup p,
 \qquad M=(N_*\setminus p)\cup T,
 \quad T\subset p.
 \label{sc:sector-support}
\end{equation}
Moreover, for arbitrary creation operators $C_1,\ldots,C_k$,
\begin{equation}
 \operatorname{ad}_{C_1}\cdots\operatorname{ad}_{C_k}(V_p)=0
 \quad\hbox{if }k>2s.
 \label{sc:depth}
\end{equation}
\end{lemma}
\begin{proof}
The commutator maps of commuting creation operators commute. A factor
whose support misses $p$ can therefore be placed next to $V_p$ and
annihilates it. Expand a surviving nested commutator into products with
some elementary operators to the left of $V_p$ and the rest to its right.
Outside $p$, the operator $V_p$ changes nothing; every site created there
must remain excited, and a site created twice makes that product zero.
This proves \eqref{sc:sector-support}. In a nonzero product the creation
supports on each side of $V_p$ are pairwise disjoint. Since each meets
$p$, at most $|p|\le s$ can occur on either side. More than $2s$
factors therefore give zero.
\end{proof}

\begin{lemma}[Sector norm estimate]
\label{sc:sector-estimate}
For $0\le k\le2s$ and creation operators $C_1,\ldots,C_k$, define
$\mathcal B_k(C_1,\ldots,C_k)$ by its coefficients
\begin{equation}
 b_M=(H_0^M)^{-1}\Pi_M
       \operatorname{ad}_{C_1}\cdots\operatorname{ad}_{C_k}(V)\Omega,
 \qquad M\ne\varnothing.
 \label{sc:Bk}
\end{equation}
For $k=0$ the commutator product is the identity. Then
\begin{equation}
 \|\mathcal B_k(C_1,\ldots,C_k)\|_1
 \le\frac{DJ}{\Delta}\beta_k(s)
                      \prod_{i=1}^k\|C_i\|_1.
 \label{sc:multilinear}
\end{equation}
\end{lemma}
\begin{proof}
Expand each $C_i$ into its elementary terms, and fix a tuple
$(M_1,\ldots,M_k,p)$. Denote its commutator by $X$ as above. Its norm
satisfies
\begin{equation}
 \|X\Omega\|\le2^kJ\prod_{i=1}^k\|c^{(i)}_{M_i}\|.
 \label{sc:tuple-norm}
\end{equation}
Only tuples with $M_i\cap p\ne\varnothing$ need be retained.
Fix a site $u$ when estimating the norm in
\eqref{sc:creation-norm}. We split the sum into tuples with $u\in p$
and those with $u\notin p$. In the latter case, a nonzero output
sector containing $u$ requires $u\in M_j$ for at least one $j$;
summing over all such $j$ is an upper bound.

If $u\in p$, put $n=|N_*\setminus p|$. A contributing output
sector has the form \eqref{sc:sector-support} with $u\in T$ and
$|M|=n+|T|$. Cauchy--Schwarz and
\eqref{sc:orthogonal-sectors} give
\begin{align}
 \sum_{M\ni u}\frac{\|\Pi_MX\Omega\|}{|M|}
 &\le\left(\sum_{\substack{T\subset p\\u\in T}}
                     \frac{1}{(n+|T|)^2}\right)^{1/2}\|X\Omega\|
 \nonumber\\
 &\le\gamma_0(s)\|X\Omega\|.
 \label{sc:type-zero}
\end{align}
For each fixed $p$ and $i$,
\begin{equation}
 \sum_{M_i:M_i\cap p\ne\varnothing}\|c^{(i)}_{M_i}\|
 \le\sum_{v\in p}\sum_{M_i\ni v}\|c^{(i)}_{M_i}\|
 \le s\|C_i\|_1.
 \label{sc:incidence-sum}
\end{equation}
There are at most $D$ choices of $p$ containing $u$. After using
$\|(H_0^M)^{-1}\|\le1/(\Delta|M|)$, the contribution of
\eqref{sc:type-zero} is bounded by
\[
 \frac{DJ}{\Delta}\,2^k\gamma_0(s)s^k
                     \prod_i\|C_i\|_1.
\]

Now suppose $u\notin p$ and $u\in M_j$. Then
$n=|N_*\setminus p|\ge1$ and $|M_j|\le n+s$. Hence
\begin{align}
 \sum_M\frac{\|\Pi_MX\Omega\|}{|M|}
 &\le\left(\sum_{t=0}^{s}\frac{\binom{s}{t}}{(n+t)^2}
          \right)^{1/2}\|X\Omega\|
 \nonumber\\
 &\le\frac{\gamma_1(s)}{|M_j|}\|X\Omega\|.
 \label{sc:type-j}
\end{align}
For the last inequality, each function
$(n+s)/(n+t)$ is nonincreasing for $n\ge1$, since $t\le s$.
Thus its square is bounded by $(s+1)^2/(t+1)^2$, and
\eqref{sc:gammas} applies. This proves the estimate for every $n$;
no finite enumeration of sector sizes is needed.

For fixed $M_j$, at most $D|M_j|$ members $p$ intersect $M_j$.
This cancels the denominator $|M_j|$ in \eqref{sc:type-j}.
Sum the other $k-1$ supports using \eqref{sc:incidence-sum},
and then sum $M_j\ni u$. For this $j$ the result is bounded by
\[
 \frac{DJ}{\Delta}\,2^k\gamma_1(s)s^{k-1}
                         \prod_i\|C_i\|_1.
\]
Adding the $k$ possible choices of $j$ proves
\eqref{sc:multilinear}. For $k=0$, only the first case occurs,
with $N_*=\varnothing$, and its bound is $DJ\gamma_0(s)/\Delta$.
\end{proof}

\begin{proof}[Proof of Theorem~\ref{sc:stability} in finite dimension]
For a creation operator $C$, put
\begin{equation}
 W(C)=e^CVe^{-C}
     =\sum_{k=0}^{2s}\frac{1}{k!}\operatorname{ad}_C^k(V).
 \label{sc:W}
\end{equation}
The equality is exact by Lemma~\ref{sc:support}.
By \eqref{sc:free-commutator}, $e^{-C}\Omega$ is an eigenvector
of $H(\varepsilon)$ precisely when
\begin{equation}
 c_M=\varepsilon(H_0^M)^{-1}\Pi_M W(C)\Omega
 \quad(M\ne\varnothing).
 \label{sc:KT}
\end{equation}
Its eigenvalue is $E(\varepsilon)=\varepsilon\langle\Omega,W(C)\Omega\rangle$.
Indeed, conjugating $H$ by $e^C$ gives
$H_0+[C,H_0]+\varepsilon W(C)$; its nonvacuum components on
$\Omega$ vanish exactly when \eqref{sc:KT} holds.

Let $F_\varepsilon$ denote the right side of \eqref{sc:KT},
assembled as a creation operator. Lemma~\ref{sc:sector-estimate} gives
\begin{equation}
 \|F_\varepsilon(C)\|_1\le\lambda q_s(\|C\|_1).
 \label{sc:F-bound}
\end{equation}
The same multilinear estimate, applied to the difference of two
$k$-fold commutators, gives, whenever $\|C\|_1,\|C'\|_1\le r$,
\begin{equation}
 \|F_\varepsilon(C)-F_\varepsilon(C')\|_1
 \le\lambda q_s'(r)\|C-C'\|_1.
 \label{sc:contraction}
\end{equation}
For example, the difference of the $k$-th terms is a sum of $k$
multilinear terms containing one $C-C'$ and $k-1$ factors bounded
by $r$. This explains the derivative of $q_s$ in the estimate.
By \eqref{sc:criterion}, the closed ball of radius $r$ is preserved
and the contraction factor is less than $1/2$. The Banach fixed-point
theorem supplies a solution $C$ with $\|C\|_1\le r$. In finite
volume the space of creation coefficients is complete in
\eqref{sc:creation-norm}. The vector $e^{-C}\Omega$ is nonzero,
since its vacuum coefficient is one.

We next prove isolation and simplicity of this eigenvalue. Suppose
$H\phi=(E+\delta)\phi$ with $|\delta|<\Delta/2$, where $\phi$
is linearly independent of $e^{-C}\Omega$. There is a unique
creation operator $B$ and scalar $b_0$ such that
\[
 e^C\phi=(B+b_0)\Omega.
\]
Here $B\ne0$. The nonvacuum components of the eigenvalue equation
after conjugation yield
\begin{equation}
 b_M=\varepsilon(H_0^M-\delta)^{-1}
                     \Pi_M[B,W(C)]\Omega.
 \label{sc:excited-equation}
\end{equation}
To verify the sign, note that
$[H_0+[C,H_0]+\varepsilon W(C),B]
 =[H_0,B]-\varepsilon[B,W(C)]$.
The creation part $[C,H_0]$ commutes with $B$.
Since $|M|\ge1$,
\begin{equation}
 \|(H_0^M-\delta)^{-1}\|
 \le\frac{1}{\Delta|M|-|\delta|}
 \le\frac{2}{\Delta|M|}.
 \label{sc:shifted-resolvent}
\end{equation}
The proof of Lemma~\ref{sc:sector-estimate} therefore applies with
an extra factor two. In \eqref{sc:W}, commutation with $B$ increases
the order by one; orders greater than $2s$ vanish. Consequently
\begin{align}
 \|B\|_1
 &\le2\lambda\sum_{k=0}^{2s-1}
       \frac{\beta_{k+1}(s)}{k!}\|C\|_1^k\|B\|_1
 \nonumber\\
 &\le2\lambda q_s'(r)\|B\|_1<\|B\|_1,
 \label{sc:gap-contradiction}
\end{align}
a contradiction. This includes $\delta=0$, so $E$ is simple.

Finally vary the real parameter from $0$ to $\varepsilon$. The same
ball and contraction factor work throughout this segment, and
\eqref{sc:contraction} implies continuous dependence of the fixed
point on the parameter. Thus $E$ is a continuous simple eigenvalue
branch, separated by at least $\Delta/2$ from every other eigenvalue.
At zero it is the ground eigenvalue of $H_0$. Continuity of the
finite-dimensional eigenvalues prevents a crossing, proving
\eqref{sc:gap}.
\end{proof}

\begin{proof}[Extension to compact-resolvent site operators]
Let $P_{u,R}=\boldsymbol{1}_{[0,R]}(h_u)$ and
$\mathcal F_R=\bigotimes_u\Ran P_{u,R}$. These are finite-dimensional
spaces containing the product vacuum. Compress $H$ to $\mathcal F_R$.
Each onsite gap remains at least $\Delta$, and each compressed
interaction is supported on the same $p$ and has norm at most $J$.
The finite-dimensional theorem therefore applies with the same constants.

In a finite system $V$ is bounded. Hence $H$ is self-adjoint on
$\Dom H_0$, bounded below, and has compact resolvent. The nested
union of the spaces $\mathcal F_R$ contains all finite linear
combinations of product eigenvectors of the $h_u$, and is an operator
core, hence also a form core, for $H_0$ and $H$.
The min--max principle gives, for each fixed eigenvalue index $j$,
\[
 E_j(H|_{\mathcal F_R})\mathbin{\downarrow}E_j(H),
\]
where $H|_{\mathcal F_R}$ denotes compression and eigenvalues are
ordered with multiplicity. If the full space has dimension at least
two, pass to the limit in
$E_1(H|_{\mathcal F_R})-E_0(H|_{\mathcal F_R})\ge\Delta/2$.
The resulting positive difference proves both the asserted gap and
simplicity. For a one-dimensional full space, the spectral exclusion
is vacuous and the ground state is simple. No norm of an unbounded
onsite operator is used in the estimates.
\end{proof}

\subsection{Explicit rational certificates}
\label{sc:certificates}

The following choices make the criterion explicit:
\begin{equation}
 \lambda_3=0.01790402,\qquad \lambda_4=0.01190829.
 \label{sc:certified-lambdas}
\end{equation}
Every terminating decimal in this subsection denotes an exact rational
number. To check the inequalities without locating a polynomial root,
use the following rational upper bounds and witnesses:
\begin{center}
\begin{tabular}{c c c c}
$s$ & $\overline\gamma_0$ & $\overline\gamma_1$ & $r$\\
\hline
3 & 1.269295517644 & 5.859465277083 & 0.03911425\\
4 & 1.464866319271 & 8.597480251020 & 0.02996979
\end{tabular}
\end{center}
The exact squares required by \eqref{sc:gammas} are
\begin{equation}
 \begin{array}{c|cc}
 s&\gamma_0(s)^2&\gamma_1(s)^2\\ \hline
 3&29/18&103/3\\
 4&103/48&887/12
 \end{array}.
 \label{sc:exact-gamma-squares}
\end{equation}
Squaring the displayed rational upper bounds verifies
$\overline\gamma_i^2\ge\gamma_i^2$.
Let $\overline q_s$ be \eqref{sc:majorant} with these upper bounds
substituted in \eqref{sc:betas}. Direct rational substitution gives,
for each of the two rows,
\begin{equation}
 r-\lambda_s\overline q_s(r)>\frac{3}{10^8},
 \qquad
 1-2\lambda_s\overline q_s'(r)>\frac{9}{10^7}.
 \label{sc:rational-slacks}
\end{equation}
These are finite rational inequalities, verifiable by clearing
denominators. Since $q_s$ and $q_s'$ have nonnegative coefficients
bounded by those of $\overline q_s$, they imply
\eqref{sc:criterion} for all $0\le\lambda\le\lambda_s$.
The certificate checks only this arithmetic; the universal operator
conclusion comes from the preceding proof. It does not follow from
finite matrix diagonalizations or from sampling interaction operators.

\subsection{Application to the specified $\SU(N)$ lattice Hamiltonian}

Let $d\ge2$ be the spatial dimension, $N\ge2$, $x=g^2>0$, and
$0\le b_N\le2N$. On a finite collection $E$ of positively oriented
cubic-lattice links, retain any subset $\mathcal P$ of the usual
elementary plaquettes whose four links belong to $E$, with no repetitions
of an anchored plaquette. Reverse
orientations denote inverses of the positive link variables.
At lattice spacing one, use
\begin{equation}
 H_{\rm lat}(x)=\frac{x}{2}\sum_{l\in E}C_l+
       \frac{b_N}{x}\sum_{p\in\mathcal P}(1-T_p),
 \qquad T_p=\frac{\operatorname{Re}\Tr U_p}{N}.
 \label{sc:KS-H}
\end{equation}
The link space is $L^2(\SU(N),\dd U)$ for normalized Haar measure;
$C_l$ is the positive group Casimir with
$\Tr(t_at_b)=\delta_{ab}/2$. The physical space is invariant under
the vertex gauge groups, with no external charges. These conventions
specify the magnetic normalization explicitly; compare the Hamiltonian
formulation in \cite{KS}. The bounds below concern the dimensionless
energies of \eqref{sc:KS-H}.

The least nonzero Casimir is
\begin{equation}
 C_F=\frac{N^2-1}{2N}\ge\frac34.
 \label{sc:CF}
\end{equation}
For completeness, write a nonzero dominant highest weight as
$\Lambda=\omega_i+\nu$, where $\nu$ is dominant and the roots have
squared length two. The Gram entries
$(\omega_i,\omega_j)=\min(i,j)(N-\max(i,j))/N$ are positive, so
$C_2(\Lambda)\ge C_2(\omega_i)$ for
$C_2(\Lambda)=(\Lambda,\Lambda+2\rho)/2$,
where $\rho=\sum_{j=1}^{N-1}\omega_j$.
Since $C_2(\omega_i)=i(N-i)(N+1)/(2N)\ge C_F$, and the
fundamental representation attains equality, Peter--Weyl decomposition
proves \eqref{sc:CF} on the entire link space.

Divide \eqref{sc:KS-H} by $s_0=xC_F/2$ and remove its scalar
magnetic constant. Precisely,
\begin{equation}
 \frac{H_{\rm lat}}{s_0}-\varepsilon|\mathcal P|I
   =\sum_{l\in E}\frac{C_l}{C_F}
        +\varepsilon\sum_{p\in\mathcal P}(-T_p),
 \qquad
 \varepsilon=\frac{2b_N}{x^2C_F}\le\frac{32}{3x^2}.
 \label{sc:normalization}
\end{equation}
The last inequality uses
$2b_N/C_F\le8N^2/(N^2-1)\le32/3$.
The normalized onsite gap is $\Delta=1$, $J=1$ because
$\|T_p\|\le1$, and each plaquette meets $s=4$ link sites.
Each link belongs to at most $D=2(d-1)$ elementary plaquettes.
These statements hold for any finite open subgraph, and also for
the usual periodic directed cubic lattices of side lengths at least
two, with positive-direction link factors distinguished.

\begin{corollary}[Explicit $N$-uniform finite-volume gauge gap]
\label{sc:gauge-corollary}
For the model \eqref{sc:KS-H}, let
\begin{equation}
 X_d^{\rm link}
  =\left(\frac{(32/3)\,2(d-1)}{0.01190829}\right)^{1/2}.
 \label{sc:X}
\end{equation}
If $x\ge X_d^{\rm link}$, the ground state is unique and gauge
invariant, and both on the full link space and on its physical
subspace,
\begin{equation}
 H_{\rm lat}-E_0\ge\frac{xC_F}{4}\ge\frac{3x}{16}
 \quad\hbox{on the orthogonal complement of the ground state}.
 \label{sc:gauge-gap}
\end{equation}
In three spatial dimensions the sufficient condition $x\ge59.9$
implies \eqref{sc:gauge-gap}, for every $N\ge2$ and
$0\le b_N\le2N$, uniformly over the stated finite volumes.
\end{corollary}
\begin{proof}
Equations \eqref{sc:normalization}--\eqref{sc:X} give
$\lambda=\varepsilon D\le\lambda_4$.
Theorem~\ref{sc:stability} and the certificate give a normalized
gap at least $1/2$. Restoring $s_0$ gives $xC_F/4$; adding back
the scalar constant changes no gap. All Casimirs and plaquette
traces commute with the finite product of vertex gauge groups.
The unique ground line therefore carries a continuous one-dimensional
representation of a product of $\SU(N)$ groups. Such a representation
is trivial, so the ground vector is gauge invariant. The physical
space reduces $H_{\rm lat}$ and contains its ground state;
restricting the excited-space inequality cannot lower the gap.
Finally, the exact rational inequality
$(59.9)^2(0.01190829)>(32/3)\,4$ verifies the stated threshold.
\end{proof}

For reference, upward-rounded sufficient thresholds using link sites
are $42.4$, $59.9$, and $73.4$ in spatial dimensions $2$, $3$, and $4$,
respectively. The exact values from \eqref{sc:X} are approximately
$42.326$, $59.858$, and $73.310$; these approximations are not used
as downward-rounded sufficient bounds.
Alternatively, group the $d$ outgoing positive links at each vertex
into one site. Each plaquette then meets three sites, and a vertex
site meets at most $3\binom d2$ plaquettes: for each coordinate pair
its three possible roles are the lower anchor and its two forward
neighbors. With $s=3$, $D=3d(d-1)/2$, and $\lambda_3$ from
\eqref{sc:certified-lambdas}, the sufficient threshold is
$\bigl((32/3)D/\lambda_3\bigr)^{1/2}$.
In $d=2$, this permits the slightly smaller upward-rounded value
$42.3$. Link grouping gives the better displayed bound in $d=3$.

\begin{remark}[Limits of the conclusion]
\label{sc:scope}
The uniformity established here is over finite volume and $N$ at
fixed $x$. The proof does not construct an infinite-volume GNS
Hamiltonian at the explicit threshold, nor does it prove a continuum
limit or a positive physical continuum mass. An infinite-volume
strong-coupling construction with a non-explicit threshold is a
separate consequence of the stability theorem in \cite{Yarotsky},
as applied in \cite{Spectral}. A limit with fixed 't~Hooft coupling
$Nx$ eventually leaves the present window as $N$ increases.
\end{remark}

\subsection{A counting correction in the cited BDL preprint}
\label{sc:BDL-correction}

This subsection concerns the exact version
arXiv:0707.1894v1 of \cite{BDL}, Appendix~A, p.~25, in the
paragraph preceding its equation~(69). The intermediate inference
$M_j\subset M\cup\{w\}$, and hence $|M_j|\le2|M|$, does not
follow from the support conditions stated there. The following
example also verifies that this is an operator issue, rather than
only a set-theoretic possibility.

Take three distinct qubit sites $u,v,w$, their vacuum $|000\rangle$,
and
\[
 A=a_u^\dagger a_v^\dagger a_w^\dagger,
 \qquad V_{vw}=X_vX_w,
 \qquad a^\dagger=|1\rangle\langle0|,
 \quad a=|0\rangle\langle1|.
\]
Here $X$ interchanges $|0\rangle$ and $|1\rangle$.
Direct application gives
\begin{equation}
 AV_{vw}|000\rangle=0,\qquad
 V_{vw}A|000\rangle=|100\rangle,\qquad
 \langle000|a_u[A,V_{vw}]|000\rangle=-1.
 \label{sc:BDL-counterexample}
\end{equation}
Thus a nonzero term has $M_j=\{u,v,w\}$ and $M=\{u\}$.
It obeys the source's support conditions but violates
$|M_j|\le2|M|$.

The valid general estimate is
\begin{equation}
 M_j\subset M\cup p,\quad |p|\le s,\quad M\ne\varnothing
 \quad\Longrightarrow\quad
 |M_j|\le|M|+s\le(s+1)|M|.
 \label{sc:corrected-count}
\end{equation}
For two-body interactions this gives $3|M|$, which is sharp by
\eqref{sc:BDL-counterexample}. It does not invalidate the conservative
constant $2^{13}$ in the source's Lemma~3. Repeating its counting
with the factor three instead of two gives, for the possible
commutator orders $1\le k\le4$, the coefficient bound
\begin{equation}
 k\cdot4\cdot3\cdot2^k2^{k-1}
        +4\cdot2^k2^k
 \le7168<8192=2^{13}.
 \label{sc:BDL-constant-survives}
\end{equation}
Higher orders vanish, and the zeroth-order estimate is smaller.
Accordingly the correction repairs this intermediate counting step
while preserving that theorem's stated conservative constant.
No priority claim is made for identifying the correction, and no
claim about revisions other than the specified source version is
needed here.

\section{Fixed-block coercivity in the actual vacuum}
\label{vac:section}

This section proves a lower bound for the energy of fluctuations in a fixed
finite set of links, conditional on all exterior links, in the actual
finite-volume ground state.  The constant is independent of the number of
exterior links.  Its dependence on the number of selected links is explicit
and does not give a uniform bound for an extensive elimination layer.
We use a Hamiltonian of Kogut--Susskind type~\cite{KS} with the normalization
specified below.  The argument is presented in full; the associated research
record is available in~\cite{Archive}.

\subsection{Finite-volume Hamiltonian and its ground-state form}
\label{vac:setup}

Let $E$ be a finite set of oriented link representatives of a finite spatial
graph, let $\mathcal P$ be a finite collection of plaquette terms, and fix
$N\geq2$, $x>0$, $a>0$, and $b\geq0$.  Put
$M=\SU(N)^E$ and equip $M$ with normalized product Haar measure $\mathrm dU$.
For Hermitian generators satisfying
$\operatorname{Tr}(t_\alpha t_\beta)=\delta_{\alpha\beta}/2$, choose the
bi-invariant metric for which the corresponding left-invariant vector fields
$X_\alpha$ are orthonormal.  On link $\ell$, write
$C_\ell=-\sum_\alpha X_{\ell,\alpha}^2=-\Delta_\ell$.  Consider
\begin{equation}
 H=\frac{x}{2}\sum_{\ell\in E}C_\ell+V,
 \qquad
 V(U)=\frac b x\sum_{p\in\mathcal P}
       \left(1-\frac{\operatorname{Re}\operatorname{Tr}U_p}{N}\right).
 \label{vac:H}
\end{equation}
Here $U_p$ is the ordered plaquette holonomy.  Each summand in parentheses
lies in $[0,2]$.  The operator $H$ is the self-adjoint bounded-potential
perturbation of the product Laplacian, with operator domain $H^2(M)$ and
form domain $H^1(M)$.  The parameter $a$ enters the physical energy through
division by $a$; it has not been absorbed into $H$.

The vertex gauge group acts by
$U_\ell\mapsto g_{s(\ell)}U_\ell g_{t(\ell)}^{-1}$.  We assume the
plaquette terms and boundary conditions preserve this action.  Physical
functions are invariant under this group.  No gauge fixing is used.

\begin{lemma}[Actual vacuum and ground-state transform]
\label{vac:groundstate}
The lowest eigenvalue $E_0$ of $H$ is simple.  Its normalized eigenfunction
$\psi$ can be chosen smooth, strictly positive, and gauge invariant.  Set
$\mathrm d\mu=\psi^2\,\mathrm dU$.  Multiplication by $\psi$ is a unitary map
from $L^2(\mu)$ to $L^2(\mathrm dU)$, and the nonnegative self-adjoint operator
\begin{equation}
 K=\psi^{-1}(H-E_0)\psi/a
 \label{vac:K}
\end{equation}
has closed quadratic form
\begin{equation}
 k[f]=\frac{x}{2a}\int_M
          \sum_{\ell,\alpha}|X_{\ell,\alpha}f|^2\,\mathrm d\mu,
 \qquad f\in H^1(M).
 \label{vac:groundform}
\end{equation}
Its kernel is the space of constant functions.  These statements also hold
after restriction to the gauge-invariant subspaces.
\end{lemma}

\begin{proof}
The compact manifold $M$ is connected, whether or not the underlying graph
is connected.  Its heat kernel is strictly positive at every positive time.
The Feynman--Kac weight for the bounded real potential $V$ is strictly
positive, so $\exp(-tH)$ is positivity improving.  Compact resolvent gives a
lowest eigenfunction.  The absolute-value inequality for the Dirichlet form
allows a nonnegative minimizer, which is strictly positive by the strong
maximum principle.  Positivity improvement implies simplicity: a second
real ground eigenfunction orthogonal to the positive one would change sign,
contradicting the strict positivity statement.  Elliptic regularity gives
smoothness.  Gauge transformations commute with $H$ and preserve positivity
and norm; uniqueness therefore makes $\psi$ gauge invariant.

For smooth complex $f$, integration by parts and $H\psi=E_0\psi$ give
\begin{equation}
 \langle\psi f,(H-E_0)\psi f\rangle_{L^2(\mathrm dU)}
   =\frac{x}{2}\int_M
       \sum_{\ell,\alpha}|X_{\ell,\alpha}f|^2\psi^2\,\mathrm dU.
 \label{vac:transformidentity}
\end{equation}
Indeed, the terms containing $|f|^2|\nabla\psi|^2$ and the mixed derivatives
combine, after integration by parts, with $V|\psi f|^2$ to give
$E_0|\psi f|^2$.  On this fixed finite graph, $\psi$ is smooth and bounded
away from zero.  Multiplication by $\psi$ and its inverse preserve $H^1(M)$,
so the identity extends by form closure.  Unitarity follows directly from
$\mathrm d\mu=\psi^2\mathrm dU$.  The zero-energy functions have zero gradient
and are constant because $M$ is connected.  Multiplication by $\psi$
intertwines the gauge actions, which proves the final assertion.
\end{proof}

In particular, all energy estimates below subtract $E_0$, the actual
interacting ground energy.  No bound on $E_0$ uniform in volume is assumed.
The equivalence of weighted and unweighted Sobolev norms just used is needed
only for each fixed graph; its constants need not be uniform in volume.

\subsection{Conditional projection and finite-block theorem}
\label{vac:block}

Fix a nonempty subset $B\subseteq E$ with $m=|B|$, and write a configuration
as $(u,y)\in\SU(N)^B\times\SU(N)^{E\setminus B}$.  Let $p_B$ be the number
of plaquette terms that involve at least one link in $B$, counting each term
once.  Define
\begin{equation}
 Z(y)=\int\psi(v,y)^2\,\mathrm dv,
 \qquad
 \mathrm d\mu_y(u)=\frac{\psi(u,y)^2}{Z(y)}\,\mathrm du,
 \qquad
 Pf(u,y)=\int f(v,y)\,\mathrm d\mu_y(v),
 \label{vac:conditional}
\end{equation}
where $\mathrm du$ and $\mathrm dv$ are normalized product Haar measures.
Thus $P=\mathbb E_\mu[\,\cdot\mid y]$ is an orthogonal projection in
$L^2(\mu)$ and $P1=1$.  Put $Q=I-P$.

The restriction of the form $k$ to
$H^1(M)\cap\operatorname{Ran}Q$ is closed and densely defined in
$\operatorname{Ran}Q$; the density assertion is verified in the proof below.
Let $D_B$ be its associated self-adjoint operator on that subspace.  Thus
$D_B$ denotes a form compression.  It does not require the unbounded
operator product $QKQ$ to be defined on an assumed common operator domain.

Write
\begin{equation}
 C_F=\frac{N^2-1}{2N},\qquad n=N^2-1,
 \qquad
 \kappa_N(s)=
 \frac{\max_{u,v\in\SU(N)}k_s(u,v)}
      {\min_{u,v\in\SU(N)}k_s(u,v)},\quad s>0,
 \label{vac:constants}
\end{equation}
where $k_s$ is the one-link heat kernel of $\exp(-sC)=\exp(s\Delta)$
relative to normalized Haar measure.  Compactness and strict positivity make
$\kappa_N(s)$ finite.

\begin{theorem}[Exterior-volume-uniform finite-block coercivity]
\label{vac:finiteblock}
For every $s>0$,
\begin{equation}
 D_B\ \geq\ d_B(s)I,
 \qquad
 d_B(s)=\frac{xC_F}{2a}\,
       \kappa_N(s)^{-2m}
       \exp\!\left(-\frac{8sbp_B}{x^2}\right).
 \label{vac:heatbound}
\end{equation}
The bound is independent of the number of links outside $B$ when
$m,p_B,N,x,b,a$ are fixed.  The projection $P$ commutes with the gauge
action, so the same bound holds on the gauge-invariant part of
$\operatorname{Ran}Q$, if that part is nonzero.  If $bp_B=0$, the sharper
bound
\begin{equation}
 D_B\geq\frac{xC_F}{2a}I
 \label{vac:freebound}
\end{equation}
holds.
\end{theorem}

\begin{proof}
Split the potential into terms that avoid and meet $B$:
\begin{equation}
 V(u,y)=V_{\mathrm{out}}(y)+V_{\mathrm{hit}}(u,y),
 \qquad
 0\leq V_{\mathrm{hit}}\leq W_B:=\frac{2bp_B}{x}.
 \label{vac:split}
\end{equation}
Let $(U_t,Y_t)$ be independent link Brownian motions, with generator
$(x/2)\Delta$ on each link.  At any time $t>0$, the eigenfunction identity
and Feynman--Kac formula give
\begin{equation}
 e^{-tE_0}\psi(u,y)=
 \mathbb E_{u,y}\!\left[
 e^{-\int_0^t V(U_r,Y_r)\,\mathrm dr}\psi(U_t,Y_t)\right].
 \label{vac:fk}
\end{equation}
Remove only $V_{\mathrm{hit}}$ from this expectation and call the result
$F_t(u,y)$.  Positivity gives
\begin{equation}
 e^{-tW_B}F_t(u,y)\leq e^{-tE_0}\psi(u,y)\leq F_t(u,y).
 \label{vac:fksandwich}
\end{equation}
Put $s=xt/2$.  Independence of the link motions and the absence of $u$ from
$V_{\mathrm{out}}$ imply
\begin{align}
 F_t(u,y)&=\int k_s^B(u,v)A_t(y,v)\,\mathrm dv,
 \label{vac:factorization}\\
 A_t(y,v)&=\mathbb E_y\!\left[
       e^{-\int_0^t V_{\mathrm{out}}(Y_r)\,\mathrm dr}
       \psi(v,Y_t)\right]>0,
 \notag
\end{align}
where $k_s^B$ is the product of the $m$ one-link kernels.  For a fixed $y$,
the positive factor $A_t(y,v)$ is the same for every initial block
configuration $u$.  Consequently, for every $u,u',y$,
\begin{equation}
 \frac{\psi(u,y)}{\psi(u',y)}
 \leq e^{tW_B}\frac{F_t(u,y)}{F_t(u',y)}
 \leq e^{tW_B}\kappa_N(s)^m=:R_B(s).
 \label{vac:oscillation}
\end{equation}
The factor $e^{-tE_0}$ cancels in this same-time ratio.  The exterior
potential remains inside the common positive expectation in
\eqref{vac:factorization}; it has not been replaced by a bound depending on
the number of exterior plaquettes.  This is the source of the asserted
exterior-volume independence.

We next recall the product Haar Poincar\'e inequality in the chosen metric:
\begin{equation}
 \operatorname{Var}_{\mathrm du}(f)
 \leq\frac1{C_F}\int\sum_{\ell\in B,\alpha}
                 |X_{\ell,\alpha}f|^2\,\mathrm du.
 \label{vac:haar}
\end{equation}
For completeness, Peter--Weyl decomposition reduces this statement to the
smallest nonzero Casimir eigenvalue of $\SU(N)$.  With long roots of squared
length $2$, the Casimir of highest weight $\lambda$ is
$(\lambda,\lambda+2\rho)/2$.  The fundamental weight $\omega_j$ has Casimir
$j(N-j)(N+1)/(2N)$.  All pairings of fundamental weights are nonnegative,
so every nonzero dominant integral weight has Casimir at least the minimum
of these fundamental values, namely $C_F$.  On a product, the eigenvalues
add; a nonconstant component therefore has eigenvalue at least $C_F$.
This proves~\eqref{vac:haar}, including for complex functions by applying
the real inequality to their real and imaginary parts.

For any positive probability density $w$ relative to product Haar with
$\max w/\min w\leq R^2$, the minimizing-constant definition of variance and
\eqref{vac:haar} give
\begin{align}
 \operatorname{Var}_{w\mathrm du}(f)
 &\leq (\max w)\operatorname{Var}_{\mathrm du}(f)\notag\\
 &\leq\frac{\max w}{C_F}\int|\nabla_Bf|^2\,\mathrm du
 \leq\frac{R^2}{C_F}\int|\nabla_Bf|^2w\,\mathrm du,
 \label{vac:densitycomparison}
\end{align}
where $|\nabla_Bf|^2=\sum_{\ell\in B,\alpha}|X_{\ell,\alpha}f|^2$.
No Haar-mean-zero condition is imposed on $f$ in this comparison.
By~\eqref{vac:oscillation}, the density of $\mu_y$ in
\eqref{vac:conditional} has maximum-to-minimum ratio at most $R_B(s)^2$.
Applying~\eqref{vac:densitycomparison} conditionally and integrating over
the exterior marginal yields
\begin{equation}
 \|Qf\|_{L^2(\mu)}^2
 \leq\frac{R_B(s)^2}{C_F}
        \int_M|\nabla_Bf|^2\,\mathrm d\mu.
 \label{vac:conditionalpoincare}
\end{equation}
If $f\in\operatorname{Ran}Q$, the exact ground-state form therefore gives
\begin{equation}
 k[f]\geq\frac{x}{2a}\int_M|\nabla_Bf|^2\,\mathrm d\mu
       \geq\frac{xC_F}{2aR_B(s)^2}\|f\|_{L^2(\mu)}^2.
 \label{vac:discardedform}
\end{equation}
Since $t=2s/x$ and $W_B=2bp_B/x$, this is precisely
\eqref{vac:heatbound}.

To justify the form restriction, observe that $P$ maps smooth functions to
smooth functions.  On each fixed compact graph it is bounded on $H^1(M)$:
differentiating its conditional integral in an exterior direction produces
the conditional derivative of $f$ and a covariance with
$2\nabla_y\log\psi$.  These coefficients are smooth and bounded, and
conditional Cauchy--Schwarz gives the required fixed-graph bound.  Derivatives
in the $B$ directions vanish for $Pf$.  Thus $Q$ is also bounded on
$H^1(M)$.  If $f\in H^1(M)\cap\operatorname{Ran}Q$, approximate $f$ in
$H^1$ by smooth $f_j$ and use $Qf_j\to f$.  This proves smooth density in
the restricted form domain.  Smooth density in $L^2$ and boundedness of $Q$
also prove that the restriction is densely defined in
$\operatorname{Ran}Q$.  It is closed because this subspace is closed in
$L^2(\mu)$ and $k$ is closed.  Inequality~\eqref{vac:discardedform} consequently
extends to the entire restricted form domain without any uniform derivative
bound on $\psi$.

For a fixed gauge transformation, the selected and exterior link coordinates
transform separately.  Both transformations preserve their Haar measures,
and $\psi$ is invariant.  Changing variables in
\eqref{vac:conditional} therefore shows that $P$ intertwines the gauge
action.  The forms and their restrictions are gauge invariant, proving the
physical-sector assertion.

Finally, if $bp_B=0$, the Hamiltonian splits as a free product Laplacian on
$B$ plus an exterior operator.  Its positive ground function is constant on
$B$.  The conditional measures are exactly product Haar, and
\eqref{vac:haar} gives~\eqref{vac:freebound}.
\end{proof}

\subsection{An explicit estimate at weak coupling}
\label{vac:explicit}

Let $D_N$ be the diameter of $\SU(N)$ in the specified Casimir metric.
The next heat-kernel estimate supplies explicit constants in
Theorem~\ref{vac:finiteblock}.

\begin{lemma}[Compact heat-kernel ratio]
\label{vac:heatlemma}
For every $s>0$,
\begin{equation}
 \kappa_N(s)\leq 2^n
       \exp\!\left(\frac{3D_N^2}{4s}\right),
 \qquad n=N^2-1.
 \label{vac:kappabound}
\end{equation}
\end{lemma}

\begin{proof}
The bi-invariant metric has nonnegative Ricci curvature.  For a positive
solution of $u_t=\Delta u$ on a compact $n$-dimensional manifold with
nonnegative Ricci curvature, the Li--Yau estimate~\cite{LiYau} is
\begin{equation}
 |\nabla\log u|^2-\partial_t\log u\leq\frac{n}{2t}.
 \label{vac:liyau}
\end{equation}
Here is its compact, boundaryless derivation.  Set $f=\log u$,
$w=|\nabla f|^2-f_t$, and $\mathcal F=tw$.  The heat equation, Bochner
identity, and $|\operatorname{Hess}f|^2\geq(\Delta f)^2/n$ imply
\begin{equation}
 (\Delta-\partial_t)\mathcal F
 \geq\frac{2\mathcal F^2}{nt}-\frac{\mathcal F}{t}
       -2\langle\nabla f,\nabla\mathcal F\rangle.
 \label{vac:bochnerheat}
\end{equation}
For smooth positive initial data, $\mathcal F=0$ at $t=0$.  A positive
space-time maximum exceeding $n/2$ contradicts~\eqref{vac:bochnerheat}
and the maximum principle.  This proves~\eqref{vac:liyau}.  For the heat
kernel, start after a positive time shift and then let that shift tend to
zero.  Compactness removes boundary and escape issues.

Along a constant-speed minimizing geodesic $\gamma$ between times
$0<t_1<t_2$, inequality~\eqref{vac:liyau} gives
\[
 \frac{\mathrm d}{\mathrm dt}f(\gamma(t),t)
 \geq-\frac{|\dot\gamma(t)|^2}{4}-\frac{n}{2t}.
\]
Integration yields
\begin{equation}
 u(z,t_1)\leq u(w,t_2)
       \left(\frac{t_2}{t_1}\right)^{n/2}
       \exp\!\left(\frac{\operatorname{dist}(z,w)^2}
                         {4(t_2-t_1)}\right).
 \label{vac:harnack}
\end{equation}
Apply this to $k_t(v,\cdot)$ between $s$ and $2s$ and integrate the later
point $w$ against normalized Haar measure.  Total heat-kernel mass one gives
\begin{equation}
 \max k_s\leq 2^{n/2}e^{D_N^2/(4s)}.
 \label{vac:heatupper}
\end{equation}
Apply~\eqref{vac:harnack} between $s/2$ and $s$ and integrate the earlier
point $z$ instead.  The result is
\begin{equation}
 \min k_s\geq2^{-n/2}e^{-D_N^2/(2s)}.
 \label{vac:heatlower}
\end{equation}
The two inequalities prove~\eqref{vac:kappabound}.  Normalizing Riemannian
volume to Haar probability multiplies the heat kernel by a constant and
does not change the logarithmic differential estimate; its mass in the
normalized measure is one, as used above.
\end{proof}

\begin{corollary}[Explicit finite-block lower bound]
\label{vac:exponential}
If $bp_B>0$, then
\begin{equation}
 D_B\geq\frac{xC_F}{2a}\,2^{-2mn}
                   e^{-C_B/x}I,
 \qquad C_B=4D_N\sqrt{3mbp_B}.
 \label{vac:explicitbound}
\end{equation}
\end{corollary}

\begin{proof}
Substituting~\eqref{vac:kappabound} into~\eqref{vac:heatbound} gives the
lower certificate
\[
 \frac{xC_F}{2a}\,2^{-2mn}
 \exp\!\left(-\frac{3mD_N^2}{2s}-\frac{8sbp_B}{x^2}\right).
\]
The positive exponent is minimized at
\[
 s_* =\frac{xD_N}{4}\sqrt{\frac{3m}{bp_B}},
 \qquad
 \frac{3mD_N^2}{2s_*}+\frac{8s_*bp_B}{x^2}
       =\frac{4D_N\sqrt{3mbp_B}}{x}.
\]
This proves~\eqref{vac:explicitbound}.  The case $bp_B=0$ is already covered
by~\eqref{vac:freebound} and does not require this optimization.
\end{proof}

The exponential dependence is a rigorous lower certificate, not an
asymptotic formula for the actual discarded spectrum as $x\downarrow0$.
In particular, the estimate is not uniform in $m$ or $N$.

\subsection{Conditional scale budget and the remaining limitation}
\label{vac:budget}

One can quantify the strength of this certificate without asserting a
renormalization theorem.  Suppose, as an additional hypothesis, that the same
fixed-block estimate with constants
$c_0=(C_F/2)2^{-2mn}$ and $C=C_B$ were available for every relevant exact
coarse operator.  Consider the illustrative scale and coupling sequences
\begin{equation}
 a_r=A L^{-r},\qquad x_r=(\sigma+qr)^{-1},\qquad r\geq1,
 \quad A>0,\ L>1,\ \sigma>0,\ q\geq0.
 \label{vac:modelscales}
\end{equation}
The reciprocal costs of the certified bounds
$d_r^{\mathrm{cert}}=c_0x_r e^{-C/x_r}/a_r$ then satisfy the exact identity
\begin{equation}
 \sum_{r\geq1}\frac1{d_r^{\mathrm{cert}}}
 =\frac{A}{c_0}e^{C\sigma}
       \sum_{r\geq1}(\sigma+qr)\rho^r,
 \qquad \rho=\frac{e^{Cq}}{L}.
 \label{vac:budgetseries}
\end{equation}
This series is finite exactly when $\rho<1$, in which case it equals
\begin{equation}
 \frac{A}{c_0}e^{C\sigma}
 \left(\frac{\sigma\rho}{1-\rho}
       +\frac{q\rho}{(1-\rho)^2}\right).
 \label{vac:budgetvalue}
\end{equation}
No actual Yang--Mills coupling trajectory, choice of scale factor, or
stability of the finite-block theorem under exact elimination is proved by
this calculation.  If $\rho\geq1$, this particular certificate does not
give a finite upper bound on the inverse-gap cost; the true discarded gaps
could still give a finite sum.

Theorem~\ref{vac:finiteblock} concerns one fixed finite block in the original
Hamiltonian~\eqref{vac:H}.  A whole elimination layer generally contains an
extensive number of links.  If that layer is used as $B$, the factor
$2^{-2mn}$ deteriorates exponentially in $m$, and $C_B$ grows proportionally
to $m$ when $p_B$ does.  Eliminating bounded blocks one after another instead
requires estimates for the successive exact coarse operators and control of
the number of elimination steps.  Neither follows from the theorem: the
generated operators have not been shown to retain the form~\eqref{vac:H},
and conditional bounds for separate blocks do not by themselves provide a
collective conditional Poincar\'e inequality in an interacting vacuum.
Thus the proved exterior-volume uniformity leaves open a whole-layer
estimate for the actual compact-link ground measure that persists under
exact coarse graining.

\section{Exact elimination and its induced inner product}\label{sec:schur}

The estimates in the preceding sections concern genuine finite-lattice
operators. To use such estimates along a sequence of scales one also needs
a comparison between the fine and coarse excitation problems. This section
proves an abstract comparison; identifying a Yang--Mills renormalization
scheme satisfying all its hypotheses is a separate problem.

\subsection{A one-step comparison}

Let $\mathcal H=P\oplus Q$ be a finite-dimensional complex Hilbert space.
Suppose $H\ge0$ has nonzero kernel and block decomposition
\begin{equation}
 H=\begin{pmatrix}A&B^*\\B&D\end{pmatrix},
 \qquad D\ge dI_Q,\quad d>0 .
\end{equation}
Energies have already been shifted by the \emph{actual} ground energy.
Define
\begin{equation}\label{eq:schur-def}
 C=D^{-1}B,\quad S=A-B^*D^{-1}B,\quad
 Jp=(p,-Cp),\quad G=J^*J=I+C^*C,\quad
 K=G^{-1/2}SG^{-1/2}.
\end{equation}
$S$ is a quadratic energy on the retained coordinates, while $G$ is the
squared norm of their energy-minimizing graph lift. The normalized map
$U=JG^{-1/2}$ is an isometry.

\begin{proposition}[Metric-normalized Schur gap comparison]\label{prop:schur}
If $\operatorname{spec}K\subset\{0\}\cup[\gamma,\infty)$ for $\gamma>0$,
then
\begin{equation}\label{eq:schur-gap}
 H\ge\left(\gamma^{-1}+d^{-1}\right)^{-1}
                (I-\Pi_{\ker H}).
\end{equation}
The kernels of $K$ and $H$ have the same dimension. If $K$ has a nonzero
excited sector, then also $\gap H\le\gap K$.
\end{proposition}

\begin{proof}
For $r=q+Cp$, completion of squares gives
\begin{equation}\label{eq:completion}
 \langle(p,q),H(p,q)\rangle=\langle p,Sp\rangle+\langle r,Dr\rangle.
\end{equation}
Both summands are nonnegative, and hence $\ker H=J\ker S$.
Decompose $p=p_0+p_\perp$ orthogonally in the $G$ inner product,
with $p_0\in\ker S$. Since $\ker K=G^{1/2}\ker S$,
$Jp_\perp\perp\ker H$ and
\[
 \|Jp_\perp\|^2\le \gamma^{-1}\langle p,Sp\rangle.
\]
The decomposition $(p,q)=Jp_0+Jp_\perp+(0,r)$ therefore implies
\begin{align*}
 \operatorname{dist}((p,q),\ker H)
 &\le \|Jp_\perp\|+\|r\|\\
 &\le \gamma^{-1/2}\langle p,Sp\rangle^{1/2}
       +d^{-1/2}\langle r,Dr\rangle^{1/2}.
\end{align*}
Squaring and applying scalar Cauchy--Schwarz proves
\eqref{eq:schur-gap}. The isometry $U$ identifies $\ker K$ with $\ker H$.
For the upper comparison, minimize the fine Rayleigh quotient over
$U((\ker K)^\perp)$; on that subspace it equals the coarse Rayleigh quotient.
\end{proof}

\begin{remark}[Closed-form extension]\label{rem:schur-domain}
The same proof applies to a closed nonnegative form on $P\oplus Q$
when $D\ge dI$ is self-adjoint, $C:P\to Q$ is bounded, $s$ is a closed
nonnegative form, and the fine form has exactly the representation
\[
 h[(p,q)]=s[p]+\|D^{1/2}(q+Cp)\|^2
\]
on the domain
\[
 \{(p,q):p\in\Dom s,\ q+Cp\in\Dom D^{1/2}\}.
\]
Here $\Dom s$ denotes the form domain. The coarse operator is the
self-adjoint operator associated with
$k[z]=s[G^{-1/2}z]$ on $G^{1/2}\Dom s$, not an unchecked product
of unbounded operators. The distance argument extends on this domain.
An application to an unbounded lattice Hamiltonian must verify these
conditions; a finite-matrix computation does not verify them.
No extension to unbounded $C$ is asserted.
\end{remark}

\subsection{Why the metric matters}

For $M\in\mathbb R$ consider
\begin{equation}
 H_M=\begin{pmatrix}
 0&0&0\\
 0&1+M^2&M\\
 0&M&1
 \end{pmatrix},\qquad
 P=\operatorname{span}\{e_1,e_2\},\quad Q=\operatorname{span}\{e_3\}.
\end{equation}
Then $D=1$, $S=\operatorname{diag}(0,1)$ and
$G=\operatorname{diag}(1,1+M^2)$, whereas
\begin{equation}
 \gap H_M=\frac{M^2+2-\sqrt{M^4+4M^2}}2\longrightarrow0
 \quad\text{as }|M|\longrightarrow\infty .
\end{equation}
Thus positive lower bounds for $D$ and the \emph{unweighted} Schur
operator $S$ give no positive uniform lower bound for the fine gap.
The correct coarse gap is $(1+M^2)^{-1}$, and
Proposition~\ref{prop:schur} yields the consistent lower bound
$(M^2+2)^{-1}$.

\subsection{An additive inverse-gap budget}

\begin{corollary}[Conditional multiscale bound]\label{cor:multiscale}
Suppose $H_0,\ldots,H_k$ form a sequence of exact metric-normalized
Schur reductions satisfying Proposition~\ref{prop:schur}, or its
closed-form extension, in the same physical energy units. At step $j$
assume the discarded block satisfies $D_j\ge d_jI$ with $d_j>0$.
If $H_k$ has gap at least $m_*>0$, then
\begin{equation}\label{eq:budget}
 \gap H_0\ge
 \left(m_*^{-1}+\sum_{j=0}^{k-1}d_j^{-1}\right)^{-1}.
\end{equation}
\end{corollary}
\begin{proof}
Iterate \eqref{eq:schur-gap} backwards from the terminal operator.
\end{proof}

This avoids a fixed multiplicative loss per scale, but every operator and
metric in the hypothesis must be the actual one generated by elimination.
For example, assume \emph{as a model profile}
$a_r=AL^{-r}$, $x_r=(\sigma+qr)^{-1}$ for $r\ge1$,
where $A,\sigma>0$, $L>1$ and $q\ge0$. If one could prove
$d_r\ge cx_r/a_r$ with one $c>0$, then
\begin{equation}
 \sum_{r\ge1}d_r^{-1}
 \le\frac{A}{c}\left(\frac{\sigma}{L-1}
                   +\frac{qL}{(L-1)^2}\right).
\end{equation}
The geometric summation is exact. Neither the assumed trajectory nor
the nonlinear discarded-sector estimate is established here.

\subsection{Quantum Schur elimination need not preserve quadratic form}

A separate obstruction is that the normalized reduction above is not
the same operation as integrating a Gaussian coordinate density.
Let $H=\mathrm d\Gamma(h)$ be the bosonic second quantization of
\[
 h=\begin{pmatrix}2&1\\1&2\end{pmatrix}.
\]
Retain the subspace spanned by all $|n,0\rangle$, so that the retained
space contains the true Fock vacuum. Total particle number is preserved
by $H$ and by all the reduction maps. In the one-particle sector,
\[
 S_1=\frac32,\qquad G_1=\frac54,\qquad K_1=\frac65 .
\]
In the two-particle sector the fine matrix is
\[
 H_2=\begin{pmatrix}
 4&\sqrt2&0\\
 \sqrt2&4&\sqrt2\\
 0&\sqrt2&4
 \end{pmatrix}
\]
in the basis $|2,0\rangle,|1,1\rangle,|0,2\rangle$. Direct inversion gives
\[
 S_2=\frac{24}{7},\qquad G_2=\frac{58}{49},
 \qquad K_2=\frac{84}{29}\ne 2K_1.
\]
Since $h\ge I$, the normalized graph restriction in the $n$-particle
sector is bounded below by $nI$. Consequently $K_1$ and $K_2<3$ are
indeed the first two positive levels. They are not equally spaced.
The full Fock-space form construction is legitimate: on the discarded
part of the $n$-particle sector, $D_n\ge nI$ and $\|B_n\|=\sqrt n$,
so $\|D_n^{-1}B_n\|\le n^{-1/2}$ and the graph operator is bounded.
This example rules out automatic closure of the oscillator class under
this quantum elimination. It does not rule out controlling the generated
operators in a larger class.

\section{The remaining Yang--Mills obligations}\label{sec:open}

The results above do not establish the hypotheses of a continuum
mass-gap theorem. The following are separate missing implications.

\begin{enumerate}
\item \textbf{Collective nonlinear coercivity.}
The finite-block bound must be replaced or supplemented by an estimate
for an entire elimination layer in the actual compact-link vacuum,
uniform in volume and ultraviolet cutoff. Its validity for each fixed
block does not imply such a collective bound.

\item \textbf{Control after exact elimination.}
The estimates must survive the generated interactions, the new vacuum,
the induced coarse metric and the appropriate form domains.
The reduced operator has not been identified with another Hamiltonian
of the bare Wilson form \eqref{eq:hamiltonian}. Substitution of a running
coupling into the strong-coupling corollary is therefore insufficient.

\item \textbf{A physical terminal gap.}
At a fixed physical coarse scale, one must control the actual terminal
operator and its metric. The explicit bound at large bare $g^2$ provides
a possible comparison target, but no comparison from the actual
renormalized operator to that target is proved here.

\item \textbf{A nontrivial continuum field theory and spectral transfer.}
One must construct compatible renormalized observables, their limiting
correlations and a field theory with the requisite positivity, locality,
covariance and Yang--Mills short-distance behavior. A positive finite
physical gap requires an estimate uniform along the relevant limit.
There must also be a nonzero finite-energy sector: exclusion of spectrum
from an interval in a vacuum-only Hilbert space would not solve the
problem posed in~\cite{JW}.
\end{enumerate}

The scale distinction is especially important. The strong-coupling
bound is in lattice units. At a fixed admissible $x$, its physical lower
bound is proportional to $x/a$ and diverges as $a\downarrow0$.
Along a weak bare-coupling continuum trajectory, $x(a)$ instead tends
to zero, outside the theorem's certified window. Neither observation
disproves a continuum mass gap; both prevent an inference of such a
gap from the fixed-coupling estimate.

There is also a concrete obstruction to treating every configuration
as a positive quadratic fluctuation problem. On an even periodic
cubic lattice in spatial dimension at least two, choose central
$\SU(2)$ links
\[
 U_i(z)=(-I)^{\sum_{k<i}z_k}.
\]
Every elementary plaquette is $-I$. For a real link field $v$, vary
$U_i(z;t)=U_i(z)\exp(itv_i(z)\sigma_3/2)$. Then
\begin{equation}\label{eq:negative-hessian}
 \left.\frac{\mathrm d^2}{\mathrm dt^2}
 \frac{b}{x}\sum_p\left(1-\frac{\operatorname{Re}\Tr U_p(t)}2\right)
 \right|_{t=0}
 =-\frac{b}{4x}\sum_p|(\mathrm dv)_p|^2 .
\end{equation}
A circulation around one square has zero divergence, zero block mean
in each link component and nonzero curl. At this central background the
adjoint transports are the identity, so these are also the covariant
conditions. It is a discarded direction for the corresponding
componentwise block average and has strictly negative magnetic
Hessian for $b>0$. Thus an all-background positive-Hessian comparison
cannot provide the required nonlinear estimate. A negative Hessian of
the potential does not imply a negative excitation Hamiltonian or a
small quantum gap. Large-field regions require their own estimates.

The infinite-volume issue is kept distinct from uniformity in finite
volume. The general stability theorem proved here supplies a common
finite-volume lower bound. We do not claim in this paper to construct
the infinite-volume state and its GNS Hamiltonian at the explicit
threshold $59.9$. The earlier application of Yarotsky's results
addresses that construction in a sufficiently strong, nonexplicit
window~\cite{Yarotsky,Spectral}.

\section{Evidence, reproducibility and attribution}\label{sec:evidence}

The supplementary source archive contains the full \LaTeX{} project,
a standard-library Python verifier, exact arithmetic results, selected
historical proof notes, and selected Lean source certificates with their
toolchain manifest. The principal analytic proofs are in this paper.
The machine checks have the following narrower scope:
\begin{itemize}
\item rational upper bounds for the square-root sector constants,
the majorant supersolution inequalities and safe coupling thresholds;
\item the finite counting correction and its three-qubit matrix element;
\item exact finite-dimensional Schur and oscillator arithmetic.
\end{itemize}
The Python verifier uses rational arithmetic for its acceptance
conditions; decimal displays are explanatory only. No random search,
floating-point eigensolver, or finite sample is used to infer an
all-volume theorem. The optional Lean certificates concern scalar
inequalities, finite sets and a finite matrix calculation. They do not
formalize the analytic stability theorem, the Feynman--Kac construction,
Li--Yau theory or continuum quantum field theory.

The archive's \path{README.md} gives build and verification commands,
and \path{MANIFEST.json} records the included files and their SHA-256
hashes. Historical notes are included for provenance and may retain
their original notation or approximate threshold displays. The
statements and safe sufficient thresholds in the present manuscript
govern this release. In particular, $73.4$, rather than a downward
rounded $73.3$, is a safe decimal upper rounding of the displayed
four-spatial-dimensional link threshold when that comparison is used.
No claim of global novelty is inferred from the absence of an identified
prior formula.

\subsection*{AI assistance and review status}

OpenAI Codex assisted with source inspection, drafting, mathematical
consistency checks, exact verification code, typesetting and release
preparation. Separate AI-assisted readings were used to check the
mathematical arguments and the scope of the claims. This is not
independent human peer review, and it does not imply that every
derivation has been manually checked by the named author. The release
includes a review record stating the performed checks and their
limitations. No AI system is listed as an author.

\subsection*{License}

This manuscript and its \LaTeX{} source are licensed under Creative
Commons Attribution 4.0 International
(\href{https://creativecommons.org/licenses/by/4.0/}{CC BY 4.0}).
The accompanying repository-derived software and formal proof sources
retain their GNU Affero General Public License, version 3 or later.
The manuscript license does not relicense that software.


\begin{thebibliography}{9}
\bibitem{BDL}
S.~Bravyi, D.~P.~DiVincenzo and D.~Loss,
\emph{Polynomial-time algorithm for simulation of weakly interacting
quantum spin systems},
Commun.\ Math.\ Phys.\ \textbf{284} (2008), 481--507.
\href{https://doi.org/10.1007/s00220-008-0574-6}{doi:10.1007/s00220-008-0574-6}.
The counting assertion examined here is in Appendix A, printed p.~25 of
\href{https://arxiv.org/abs/0707.1894v1}{arXiv:0707.1894v1}.

\bibitem{Yarotsky}
D.~A.~Yarotsky,
\emph{Quasi-particles in weak perturbations of non-interacting quantum
lattice systems},
\href{https://arxiv.org/abs/math-ph/0411042v1}{arXiv:math-ph/0411042v1} (2004).

\bibitem{KS}
J.~Kogut and L.~Susskind,
\emph{Hamiltonian formulation of Wilson's lattice gauge theories},
Phys.\ Rev.\ D \textbf{11} (1975), 395--408.
\href{https://doi.org/10.1103/PhysRevD.11.395}{doi:10.1103/PhysRevD.11.395}.

\bibitem{LiYau}
P.~Li and S.-T.~Yau,
\emph{On the parabolic kernel of the Schr\"odinger operator},
Acta Math.\ \textbf{156} (1986), 153--201, Theorem~1.1.
\href{https://doi.org/10.1007/BF02399203}{doi:10.1007/BF02399203}.

\bibitem{JW}
A.~Jaffe and E.~Witten,
\emph{Quantum Yang--Mills theory},
Clay Mathematics Institute problem description, Section~4.
\url{https://www.claymath.org/wp-content/uploads/2022/06/yangmills.pdf}.

\bibitem{Spectral}
C.~P.~Zimmerman,
\emph{From Scalar Certificates to Spectral Bounds: Radial Pulsations,
Discrete Hardy Forms, and SU(N) Lattice Gauge Theory}, version~1.0,
September 22, 2026.
\href{https://doi.org/10.5281/zenodo.22895866}{doi:10.5281/zenodo.22895866}.

\bibitem{Archive}
C.~P.~Zimmerman,
\emph{The Zimmerman Programme}, software and research archive,
version v9.0.0, October 5, 2026.
\href{https://doi.org/10.5281/zenodo.23169118}{doi:10.5281/zenodo.23169118}.
\end{thebibliography}
\end{document}

% SOURCE sections/open_and_evidence.tex
\section{The remaining Yang--Mills obligations}\label{sec:open}

The results above do not establish the hypotheses of a continuum
mass-gap theorem. The following are separate missing implications.

\begin{enumerate}
\item \textbf{Collective nonlinear coercivity.}
The finite-block bound must be replaced or supplemented by an estimate
for an entire elimination layer in the actual compact-link vacuum,
uniform in volume and ultraviolet cutoff. Its validity for each fixed
block does not imply such a collective bound.

\item \textbf{Control after exact elimination.}
The estimates must survive the generated interactions, the new vacuum,
the induced coarse metric and the appropriate form domains.
The reduced operator has not been identified with another Hamiltonian
of the bare Wilson form \eqref{eq:hamiltonian}. Substitution of a running
coupling into the strong-coupling corollary is therefore insufficient.

\item \textbf{A physical terminal gap.}
At a fixed physical coarse scale, one must control the actual terminal
operator and its metric. The explicit bound at large bare $g^2$ provides
a possible comparison target, but no comparison from the actual
renormalized operator to that target is proved here.

\item \textbf{A nontrivial continuum field theory and spectral transfer.}
One must construct compatible renormalized observables, their limiting
correlations and a field theory with the requisite positivity, locality,
covariance and Yang--Mills short-distance behavior. A positive finite
physical gap requires an estimate uniform along the relevant limit.
There must also be a nonzero finite-energy sector: exclusion of spectrum
from an interval in a vacuum-only Hilbert space would not solve the
problem posed in~\cite{JW}.
\end{enumerate}

The scale distinction is especially important. The strong-coupling
bound is in lattice units. At a fixed admissible $x$, its physical lower
bound is proportional to $x/a$ and diverges as $a\downarrow0$.
Along a weak bare-coupling continuum trajectory, $x(a)$ instead tends
to zero, outside the theorem's certified window. Neither observation
disproves a continuum mass gap; both prevent an inference of such a
gap from the fixed-coupling estimate.

There is also a concrete obstruction to treating every configuration
as a positive quadratic fluctuation problem. On an even periodic
cubic lattice in spatial dimension at least two, choose central
$\SU(2)$ links
\[
 U_i(z)=(-I)^{\sum_{k<i}z_k}.
\]
Every elementary plaquette is $-I$. For a real link field $v$, vary
$U_i(z;t)=U_i(z)\exp(itv_i(z)\sigma_3/2)$. Then
\begin{equation}\label{eq:negative-hessian}
 \left.\frac{\mathrm d^2}{\mathrm dt^2}
 \frac{b}{x}\sum_p\left(1-\frac{\operatorname{Re}\Tr U_p(t)}2\right)
 \right|_{t=0}
 =-\frac{b}{4x}\sum_p|(\mathrm dv)_p|^2 .
\end{equation}
A circulation around one square has zero divergence, zero block mean
in each link component and nonzero curl. At this central background the
adjoint transports are the identity, so these are also the covariant
conditions. It is a discarded direction for the corresponding
componentwise block average and has strictly negative magnetic
Hessian for $b>0$. Thus an all-background positive-Hessian comparison
cannot provide the required nonlinear estimate. A negative Hessian of
the potential does not imply a negative excitation Hamiltonian or a
small quantum gap. Large-field regions require their own estimates.

The infinite-volume issue is kept distinct from uniformity in finite
volume. The general stability theorem proved here supplies a common
finite-volume lower bound. We do not claim in this paper to construct
the infinite-volume state and its GNS Hamiltonian at the explicit
threshold $59.9$. The earlier application of Yarotsky's results
addresses that construction in a sufficiently strong, nonexplicit
window~\cite{Yarotsky,Spectral}.

\section{Evidence, reproducibility and attribution}\label{sec:evidence}

The supplementary source archive contains the full \LaTeX{} project,
a standard-library Python verifier, exact arithmetic results, selected
historical proof notes, and selected Lean source certificates with their
toolchain manifest. The principal analytic proofs are in this paper.
The machine checks have the following narrower scope:
\begin{itemize}
\item rational upper bounds for the square-root sector constants,
the majorant supersolution inequalities and safe coupling thresholds;
\item the finite counting correction and its three-qubit matrix element;
\item exact finite-dimensional Schur and oscillator arithmetic.
\end{itemize}
The Python verifier uses rational arithmetic for its acceptance
conditions; decimal displays are explanatory only. No random search,
floating-point eigensolver, or finite sample is used to infer an
all-volume theorem. The optional Lean certificates concern scalar
inequalities, finite sets and a finite matrix calculation. They do not
formalize the analytic stability theorem, the Feynman--Kac construction,
Li--Yau theory or continuum quantum field theory.

The archive's \path{README.md} gives build and verification commands,
and \path{MANIFEST.json} records the included files and their SHA-256
hashes. Historical notes are included for provenance and may retain
their original notation or approximate threshold displays. The
statements and safe sufficient thresholds in the present manuscript
govern this release. In particular, $73.4$, rather than a downward
rounded $73.3$, is a safe decimal upper rounding of the displayed
four-spatial-dimensional link threshold when that comparison is used.
No claim of global novelty is inferred from the absence of an identified
prior formula.

\subsection*{AI assistance and review status}

OpenAI Codex assisted with source inspection, drafting, mathematical
consistency checks, exact verification code, typesetting and release
preparation. Separate AI-assisted readings were used to check the
mathematical arguments and the scope of the claims. This is not
independent human peer review, and it does not imply that every
derivation has been manually checked by the named author. The release
includes a review record stating the performed checks and their
limitations. No AI system is listed as an author.

\subsection*{License}

This manuscript and its \LaTeX{} source are licensed under Creative
Commons Attribution 4.0 International
(\href{https://creativecommons.org/licenses/by/4.0/}{CC BY 4.0}).
The accompanying repository-derived software and formal proof sources
retain their GNU Affero General Public License, version 3 or later.
The manuscript license does not relicense that software.

% SOURCE sections/schur.tex
\section{Exact elimination and its induced inner product}\label{sec:schur}

The estimates in the preceding sections concern genuine finite-lattice
operators. To use such estimates along a sequence of scales one also needs
a comparison between the fine and coarse excitation problems. This section
proves an abstract comparison; identifying a Yang--Mills renormalization
scheme satisfying all its hypotheses is a separate problem.

\subsection{A one-step comparison}

Let $\mathcal H=P\oplus Q$ be a finite-dimensional complex Hilbert space.
Suppose $H\ge0$ has nonzero kernel and block decomposition
\begin{equation}
 H=\begin{pmatrix}A&B^*\\B&D\end{pmatrix},
 \qquad D\ge dI_Q,\quad d>0 .
\end{equation}
Energies have already been shifted by the \emph{actual} ground energy.
Define
\begin{equation}\label{eq:schur-def}
 C=D^{-1}B,\quad S=A-B^*D^{-1}B,\quad
 Jp=(p,-Cp),\quad G=J^*J=I+C^*C,\quad
 K=G^{-1/2}SG^{-1/2}.
\end{equation}
$S$ is a quadratic energy on the retained coordinates, while $G$ is the
squared norm of their energy-minimizing graph lift. The normalized map
$U=JG^{-1/2}$ is an isometry.

\begin{proposition}[Metric-normalized Schur gap comparison]\label{prop:schur}
If $\operatorname{spec}K\subset\{0\}\cup[\gamma,\infty)$ for $\gamma>0$,
then
\begin{equation}\label{eq:schur-gap}
 H\ge\left(\gamma^{-1}+d^{-1}\right)^{-1}
                (I-\Pi_{\ker H}).
\end{equation}
The kernels of $K$ and $H$ have the same dimension. If $K$ has a nonzero
excited sector, then also $\gap H\le\gap K$.
\end{proposition}

\begin{proof}
For $r=q+Cp$, completion of squares gives
\begin{equation}\label{eq:completion}
 \langle(p,q),H(p,q)\rangle=\langle p,Sp\rangle+\langle r,Dr\rangle.
\end{equation}
Both summands are nonnegative, and hence $\ker H=J\ker S$.
Decompose $p=p_0+p_\perp$ orthogonally in the $G$ inner product,
with $p_0\in\ker S$. Since $\ker K=G^{1/2}\ker S$,
$Jp_\perp\perp\ker H$ and
\[
 \|Jp_\perp\|^2\le \gamma^{-1}\langle p,Sp\rangle.
\]
The decomposition $(p,q)=Jp_0+Jp_\perp+(0,r)$ therefore implies
\begin{align*}
 \operatorname{dist}((p,q),\ker H)
 &\le \|Jp_\perp\|+\|r\|\\
 &\le \gamma^{-1/2}\langle p,Sp\rangle^{1/2}
       +d^{-1/2}\langle r,Dr\rangle^{1/2}.
\end{align*}
Squaring and applying scalar Cauchy--Schwarz proves
\eqref{eq:schur-gap}. The isometry $U$ identifies $\ker K$ with $\ker H$.
For the upper comparison, minimize the fine Rayleigh quotient over
$U((\ker K)^\perp)$; on that subspace it equals the coarse Rayleigh quotient.
\end{proof}

\begin{remark}[Closed-form extension]\label{rem:schur-domain}
The same proof applies to a closed nonnegative form on $P\oplus Q$
when $D\ge dI$ is self-adjoint, $C:P\to Q$ is bounded, $s$ is a closed
nonnegative form, and the fine form has exactly the representation
\[
 h[(p,q)]=s[p]+\|D^{1/2}(q+Cp)\|^2
\]
on the domain
\[
 \{(p,q):p\in\Dom s,\ q+Cp\in\Dom D^{1/2}\}.
\]
Here $\Dom s$ denotes the form domain. The coarse operator is the
self-adjoint operator associated with
$k[z]=s[G^{-1/2}z]$ on $G^{1/2}\Dom s$, not an unchecked product
of unbounded operators. The distance argument extends on this domain.
An application to an unbounded lattice Hamiltonian must verify these
conditions; a finite-matrix computation does not verify them.
No extension to unbounded $C$ is asserted.
\end{remark}

\subsection{Why the metric matters}

For $M\in\mathbb R$ consider
\begin{equation}
 H_M=\begin{pmatrix}
 0&0&0\\
 0&1+M^2&M\\
 0&M&1
 \end{pmatrix},\qquad
 P=\operatorname{span}\{e_1,e_2\},\quad Q=\operatorname{span}\{e_3\}.
\end{equation}
Then $D=1$, $S=\operatorname{diag}(0,1)$ and
$G=\operatorname{diag}(1,1+M^2)$, whereas
\begin{equation}
 \gap H_M=\frac{M^2+2-\sqrt{M^4+4M^2}}2\longrightarrow0
 \quad\text{as }|M|\longrightarrow\infty .
\end{equation}
Thus positive lower bounds for $D$ and the \emph{unweighted} Schur
operator $S$ give no positive uniform lower bound for the fine gap.
The correct coarse gap is $(1+M^2)^{-1}$, and
Proposition~\ref{prop:schur} yields the consistent lower bound
$(M^2+2)^{-1}$.

\subsection{An additive inverse-gap budget}

\begin{corollary}[Conditional multiscale bound]\label{cor:multiscale}
Suppose $H_0,\ldots,H_k$ form a sequence of exact metric-normalized
Schur reductions satisfying Proposition~\ref{prop:schur}, or its
closed-form extension, in the same physical energy units. At step $j$
assume the discarded block satisfies $D_j\ge d_jI$ with $d_j>0$.
If $H_k$ has gap at least $m_*>0$, then
\begin{equation}\label{eq:budget}
 \gap H_0\ge
 \left(m_*^{-1}+\sum_{j=0}^{k-1}d_j^{-1}\right)^{-1}.
\end{equation}
\end{corollary}
\begin{proof}
Iterate \eqref{eq:schur-gap} backwards from the terminal operator.
\end{proof}

This avoids a fixed multiplicative loss per scale, but every operator and
metric in the hypothesis must be the actual one generated by elimination.
For example, assume \emph{as a model profile}
$a_r=AL^{-r}$, $x_r=(\sigma+qr)^{-1}$ for $r\ge1$,
where $A,\sigma>0$, $L>1$ and $q\ge0$. If one could prove
$d_r\ge cx_r/a_r$ with one $c>0$, then
\begin{equation}
 \sum_{r\ge1}d_r^{-1}
 \le\frac{A}{c}\left(\frac{\sigma}{L-1}
                   +\frac{qL}{(L-1)^2}\right).
\end{equation}
The geometric summation is exact. Neither the assumed trajectory nor
the nonlinear discarded-sector estimate is established here.

\subsection{Quantum Schur elimination need not preserve quadratic form}

A separate obstruction is that the normalized reduction above is not
the same operation as integrating a Gaussian coordinate density.
Let $H=\mathrm d\Gamma(h)$ be the bosonic second quantization of
\[
 h=\begin{pmatrix}2&1\\1&2\end{pmatrix}.
\]
Retain the subspace spanned by all $|n,0\rangle$, so that the retained
space contains the true Fock vacuum. Total particle number is preserved
by $H$ and by all the reduction maps. In the one-particle sector,
\[
 S_1=\frac32,\qquad G_1=\frac54,\qquad K_1=\frac65 .
\]
In the two-particle sector the fine matrix is
\[
 H_2=\begin{pmatrix}
 4&\sqrt2&0\\
 \sqrt2&4&\sqrt2\\
 0&\sqrt2&4
 \end{pmatrix}
\]
in the basis $|2,0\rangle,|1,1\rangle,|0,2\rangle$. Direct inversion gives
\[
 S_2=\frac{24}{7},\qquad G_2=\frac{58}{49},
 \qquad K_2=\frac{84}{29}\ne 2K_1.
\]
Since $h\ge I$, the normalized graph restriction in the $n$-particle
sector is bounded below by $nI$. Consequently $K_1$ and $K_2<3$ are
indeed the first two positive levels. They are not equally spaced.
The full Fock-space form construction is legitimate: on the discarded
part of the $n$-particle sector, $D_n\ge nI$ and $\|B_n\|=\sqrt n$,
so $\|D_n^{-1}B_n\|\le n^{-1/2}$ and the graph operator is bounded.
This example rules out automatic closure of the oscillator class under
this quantum elimination. It does not rule out controlling the generated
operators in a larger class.

% SOURCE sections/strong_coupling.tex
\section{An explicit finite-volume strong-coupling bound}
\label{sc:section}

This section gives a quantitative version of stability of a gapped product
Hamiltonian under bounded local perturbations. The argument follows the
Kirkwood--Thomas method used by Bravyi, DiVincenzo, and Loss
\cite{BDL}, with Hilbert-space norms on excitation sectors in place of
absolute values of qubit coefficients. The existence of a sufficiently
small perturbative window, including for infinite-dimensional site spaces,
is established in the literature \cite{Yarotsky}. The purpose here is to
track constants for the specified Hamiltonian lattice gauge model.
Neither the method nor the existence of strong-coupling lattice gaps is
claimed to be new. The derivation and associated records originated in
the research archive \cite{Archive}; the earlier spectral-bounds paper
\cite{Spectral} gives the corresponding non-explicit stability application.

\subsection{A finite-system stability criterion}

Let $\mathcal L$ be a nonempty finite set of sites. At each site $u$, let
$\mathcal H_u$ be a separable Hilbert space and let $h_u$ be a nonnegative
self-adjoint operator with compact resolvent. Assume that $h_u$ has a
normalized simple zero eigenvector $\Omega_u$ and that, for a common
$\Delta>0$,
\begin{equation}
 h_u\big|_{\Omega_u^\perp}\ge\Delta.
 \label{sc:onsite-gap}
\end{equation}
The inequality is vacuous if $\mathcal H_u$ is one-dimensional. On
$\mathcal H=\bigotimes_{u\in\mathcal L}\mathcal H_u$, set
\begin{equation}
 H_0=\sum_{u\in\mathcal L}h_u,\qquad
 V=\sum_{p\in\mathcal P}V_p,\qquad
 H(\varepsilon)=H_0+\varepsilon V.
 \label{sc:general-H}
\end{equation}
Here $\mathcal P$ is a finite indexed family of nonempty site sets, each
of cardinality at most an integer $s\ge1$. Each $V_p$ is self-adjoint, supported on
$p$, and has norm at most $J$. Every site belongs to at most $D$
members of $\mathcal P$, counted with their multiplicities. We may take
$D,J>0$; when there is no perturbation the conclusion below is immediate.
The coupling $\varepsilon$ is real.

Define the dimensionless parameter and positive constants
\begin{align}
 \lambda&=\frac{|\varepsilon|DJ}{\Delta},
 \label{sc:lambda}\\
 \gamma_0(s)^2&=\sum_{t=1}^{s}\frac{\binom{s-1}{t-1}}{t^2},
 &\gamma_1(s)^2&=(s+1)^2\sum_{t=0}^{s}
                        \frac{\binom{s}{t}}{(t+1)^2},
 \label{sc:gammas}\\
 \beta_0(s)&=\gamma_0(s),
 &\beta_k(s)&=2^k\bigl(\gamma_0(s)s^k+
                         k\gamma_1(s)s^{k-1}\bigr)
                         \quad(1\le k\le2s),
 \label{sc:betas}\\
 q_s(r)&=\sum_{k=0}^{2s}\frac{\beta_k(s)}{k!}r^k.
 \label{sc:majorant}
\end{align}

\begin{theorem}[Dimension-independent finite-volume stability]
\label{sc:stability}
Suppose there is a number $r>0$ such that
\begin{equation}
 \lambda q_s(r)\le r,
 \qquad 2\lambda q_s'(r)<1.
 \label{sc:criterion}
\end{equation}
Then $H(\varepsilon)$ has a simple ground eigenvalue $E_0$ and
\begin{equation}
 \operatorname{spec}H(\varepsilon)
 \subset\{E_0\}\cup[E_0+\Delta/2,\infty).
 \label{sc:gap}
\end{equation}
The criterion depends on $s,D,J,\Delta$ and $\varepsilon$, but not on
the number of sites, their Hilbert-space dimensions, or the upper
onsite spectrum.
\end{theorem}

We first prove the theorem when all site spaces are finite-dimensional.
The compact-resolvent extension is proved at the end of this subsection.

For $M\subset\mathcal L$, let
$\mathcal K_M=\bigotimes_{u\in M}\Omega_u^\perp$ and
$\Omega_M=\bigotimes_{u\in M}\Omega_u$. The full Hilbert space is the
orthogonal sum of the sectors
$\mathcal K_M\otimes\Omega_{\mathcal L\setminus M}$.
We identify a sector with $\mathcal K_M$, and write $\Pi_M$ for the
corresponding projection followed by this identification. In particular,
\begin{equation}
 \sum_{M\subset\mathcal L}\|\Pi_M\xi\|^2=\|\xi\|^2.
 \label{sc:orthogonal-sectors}
\end{equation}
Let $\Omega=\bigotimes_{u\in\mathcal L}\Omega_u$.
For $c_M\in\mathcal K_M$, $M\ne\varnothing$, define
\begin{equation}
 A[c_M]=|c_M\rangle\langle\Omega_M|\otimes
                  I_{\mathcal L\setminus M},
 \qquad
 C=\sum_{M\ne\varnothing}A[c_M],
 \qquad
 \|C\|_1=\max_{u\in\mathcal L}
                 \sum_{M\ni u}\|c_M\|.
 \label{sc:creation-norm}
\end{equation}
The decomposition of $C$ is unique, as is seen by applying it to $\Omega$.
We call $C$ a creation operator. The coefficients $c_M$ are vectors,
including when the excitations within a sector are entangled.

Two such elementary operators commute. If their supports intersect,
their product in either order is zero: a vacuum bra on an overlapping
site annihilates an excited factor. If their supports are disjoint,
they commute by tensor factorization. Consequently all creation operators
commute, and every product of more than $|\mathcal L|$ elementary creation
operators vanishes. Their exponentials are therefore finite polynomials.
Also, $\|A[c_M]\|=\|c_M\|$.

Put $H_0^M=\sum_{u\in M}h_u$ on $\mathcal K_M$. Then
\begin{equation}
 H_0^M\ge\Delta|M|,
 \qquad [H_0,A[c_M]]=A[H_0^M c_M].
 \label{sc:free-commutator}
\end{equation}
Thus a second commutator of $H_0$ with creation operators is zero.
We use $\operatorname{ad}_C(X)=[C,X]$.

\begin{lemma}[Support and commutator depth]
\label{sc:support}
Fix $p\in\mathcal P$ and elementary creation operators
$A[c_{M_1}],\ldots,A[c_{M_k}]$. Set $N_*=\bigcup_{i=1}^kM_i$ and
\[
 X=\operatorname{ad}_{A[c_{M_1}]}\cdots
                   \operatorname{ad}_{A[c_{M_k}]}(V_p).
\]
If $X\ne0$, each $M_i$ intersects $p$. If $\Pi_MX\Omega\ne0$, then
\begin{equation}
 N_*\setminus p\subset M\subset N_*\cup p,
 \qquad M=(N_*\setminus p)\cup T,
 \quad T\subset p.
 \label{sc:sector-support}
\end{equation}
Moreover, for arbitrary creation operators $C_1,\ldots,C_k$,
\begin{equation}
 \operatorname{ad}_{C_1}\cdots\operatorname{ad}_{C_k}(V_p)=0
 \quad\hbox{if }k>2s.
 \label{sc:depth}
\end{equation}
\end{lemma}
\begin{proof}
The commutator maps of commuting creation operators commute. A factor
whose support misses $p$ can therefore be placed next to $V_p$ and
annihilates it. Expand a surviving nested commutator into products with
some elementary operators to the left of $V_p$ and the rest to its right.
Outside $p$, the operator $V_p$ changes nothing; every site created there
must remain excited, and a site created twice makes that product zero.
This proves \eqref{sc:sector-support}. In a nonzero product the creation
supports on each side of $V_p$ are pairwise disjoint. Since each meets
$p$, at most $|p|\le s$ can occur on either side. More than $2s$
factors therefore give zero.
\end{proof}

\begin{lemma}[Sector norm estimate]
\label{sc:sector-estimate}
For $0\le k\le2s$ and creation operators $C_1,\ldots,C_k$, define
$\mathcal B_k(C_1,\ldots,C_k)$ by its coefficients
\begin{equation}
 b_M=(H_0^M)^{-1}\Pi_M
       \operatorname{ad}_{C_1}\cdots\operatorname{ad}_{C_k}(V)\Omega,
 \qquad M\ne\varnothing.
 \label{sc:Bk}
\end{equation}
For $k=0$ the commutator product is the identity. Then
\begin{equation}
 \|\mathcal B_k(C_1,\ldots,C_k)\|_1
 \le\frac{DJ}{\Delta}\beta_k(s)
                      \prod_{i=1}^k\|C_i\|_1.
 \label{sc:multilinear}
\end{equation}
\end{lemma}
\begin{proof}
Expand each $C_i$ into its elementary terms, and fix a tuple
$(M_1,\ldots,M_k,p)$. Denote its commutator by $X$ as above. Its norm
satisfies
\begin{equation}
 \|X\Omega\|\le2^kJ\prod_{i=1}^k\|c^{(i)}_{M_i}\|.
 \label{sc:tuple-norm}
\end{equation}
Only tuples with $M_i\cap p\ne\varnothing$ need be retained.
Fix a site $u$ when estimating the norm in
\eqref{sc:creation-norm}. We split the sum into tuples with $u\in p$
and those with $u\notin p$. In the latter case, a nonzero output
sector containing $u$ requires $u\in M_j$ for at least one $j$;
summing over all such $j$ is an upper bound.

If $u\in p$, put $n=|N_*\setminus p|$. A contributing output
sector has the form \eqref{sc:sector-support} with $u\in T$ and
$|M|=n+|T|$. Cauchy--Schwarz and
\eqref{sc:orthogonal-sectors} give
\begin{align}
 \sum_{M\ni u}\frac{\|\Pi_MX\Omega\|}{|M|}
 &\le\left(\sum_{\substack{T\subset p\\u\in T}}
                     \frac{1}{(n+|T|)^2}\right)^{1/2}\|X\Omega\|
 \nonumber\\
 &\le\gamma_0(s)\|X\Omega\|.
 \label{sc:type-zero}
\end{align}
For each fixed $p$ and $i$,
\begin{equation}
 \sum_{M_i:M_i\cap p\ne\varnothing}\|c^{(i)}_{M_i}\|
 \le\sum_{v\in p}\sum_{M_i\ni v}\|c^{(i)}_{M_i}\|
 \le s\|C_i\|_1.
 \label{sc:incidence-sum}
\end{equation}
There are at most $D$ choices of $p$ containing $u$. After using
$\|(H_0^M)^{-1}\|\le1/(\Delta|M|)$, the contribution of
\eqref{sc:type-zero} is bounded by
\[
 \frac{DJ}{\Delta}\,2^k\gamma_0(s)s^k
                     \prod_i\|C_i\|_1.
\]

Now suppose $u\notin p$ and $u\in M_j$. Then
$n=|N_*\setminus p|\ge1$ and $|M_j|\le n+s$. Hence
\begin{align}
 \sum_M\frac{\|\Pi_MX\Omega\|}{|M|}
 &\le\left(\sum_{t=0}^{s}\frac{\binom{s}{t}}{(n+t)^2}
          \right)^{1/2}\|X\Omega\|
 \nonumber\\
 &\le\frac{\gamma_1(s)}{|M_j|}\|X\Omega\|.
 \label{sc:type-j}
\end{align}
For the last inequality, each function
$(n+s)/(n+t)$ is nonincreasing for $n\ge1$, since $t\le s$.
Thus its square is bounded by $(s+1)^2/(t+1)^2$, and
\eqref{sc:gammas} applies. This proves the estimate for every $n$;
no finite enumeration of sector sizes is needed.

For fixed $M_j$, at most $D|M_j|$ members $p$ intersect $M_j$.
This cancels the denominator $|M_j|$ in \eqref{sc:type-j}.
Sum the other $k-1$ supports using \eqref{sc:incidence-sum},
and then sum $M_j\ni u$. For this $j$ the result is bounded by
\[
 \frac{DJ}{\Delta}\,2^k\gamma_1(s)s^{k-1}
                         \prod_i\|C_i\|_1.
\]
Adding the $k$ possible choices of $j$ proves
\eqref{sc:multilinear}. For $k=0$, only the first case occurs,
with $N_*=\varnothing$, and its bound is $DJ\gamma_0(s)/\Delta$.
\end{proof}

\begin{proof}[Proof of Theorem~\ref{sc:stability} in finite dimension]
For a creation operator $C$, put
\begin{equation}
 W(C)=e^CVe^{-C}
     =\sum_{k=0}^{2s}\frac{1}{k!}\operatorname{ad}_C^k(V).
 \label{sc:W}
\end{equation}
The equality is exact by Lemma~\ref{sc:support}.
By \eqref{sc:free-commutator}, $e^{-C}\Omega$ is an eigenvector
of $H(\varepsilon)$ precisely when
\begin{equation}
 c_M=\varepsilon(H_0^M)^{-1}\Pi_M W(C)\Omega
 \quad(M\ne\varnothing).
 \label{sc:KT}
\end{equation}
Its eigenvalue is $E(\varepsilon)=\varepsilon\langle\Omega,W(C)\Omega\rangle$.
Indeed, conjugating $H$ by $e^C$ gives
$H_0+[C,H_0]+\varepsilon W(C)$; its nonvacuum components on
$\Omega$ vanish exactly when \eqref{sc:KT} holds.

Let $F_\varepsilon$ denote the right side of \eqref{sc:KT},
assembled as a creation operator. Lemma~\ref{sc:sector-estimate} gives
\begin{equation}
 \|F_\varepsilon(C)\|_1\le\lambda q_s(\|C\|_1).
 \label{sc:F-bound}
\end{equation}
The same multilinear estimate, applied to the difference of two
$k$-fold commutators, gives, whenever $\|C\|_1,\|C'\|_1\le r$,
\begin{equation}
 \|F_\varepsilon(C)-F_\varepsilon(C')\|_1
 \le\lambda q_s'(r)\|C-C'\|_1.
 \label{sc:contraction}
\end{equation}
For example, the difference of the $k$-th terms is a sum of $k$
multilinear terms containing one $C-C'$ and $k-1$ factors bounded
by $r$. This explains the derivative of $q_s$ in the estimate.
By \eqref{sc:criterion}, the closed ball of radius $r$ is preserved
and the contraction factor is less than $1/2$. The Banach fixed-point
theorem supplies a solution $C$ with $\|C\|_1\le r$. In finite
volume the space of creation coefficients is complete in
\eqref{sc:creation-norm}. The vector $e^{-C}\Omega$ is nonzero,
since its vacuum coefficient is one.

We next prove isolation and simplicity of this eigenvalue. Suppose
$H\phi=(E+\delta)\phi$ with $|\delta|<\Delta/2$, where $\phi$
is linearly independent of $e^{-C}\Omega$. There is a unique
creation operator $B$ and scalar $b_0$ such that
\[
 e^C\phi=(B+b_0)\Omega.
\]
Here $B\ne0$. The nonvacuum components of the eigenvalue equation
after conjugation yield
\begin{equation}
 b_M=\varepsilon(H_0^M-\delta)^{-1}
                     \Pi_M[B,W(C)]\Omega.
 \label{sc:excited-equation}
\end{equation}
To verify the sign, note that
$[H_0+[C,H_0]+\varepsilon W(C),B]
 =[H_0,B]-\varepsilon[B,W(C)]$.
The creation part $[C,H_0]$ commutes with $B$.
Since $|M|\ge1$,
\begin{equation}
 \|(H_0^M-\delta)^{-1}\|
 \le\frac{1}{\Delta|M|-|\delta|}
 \le\frac{2}{\Delta|M|}.
 \label{sc:shifted-resolvent}
\end{equation}
The proof of Lemma~\ref{sc:sector-estimate} therefore applies with
an extra factor two. In \eqref{sc:W}, commutation with $B$ increases
the order by one; orders greater than $2s$ vanish. Consequently
\begin{align}
 \|B\|_1
 &\le2\lambda\sum_{k=0}^{2s-1}
       \frac{\beta_{k+1}(s)}{k!}\|C\|_1^k\|B\|_1
 \nonumber\\
 &\le2\lambda q_s'(r)\|B\|_1<\|B\|_1,
 \label{sc:gap-contradiction}
\end{align}
a contradiction. This includes $\delta=0$, so $E$ is simple.

Finally vary the real parameter from $0$ to $\varepsilon$. The same
ball and contraction factor work throughout this segment, and
\eqref{sc:contraction} implies continuous dependence of the fixed
point on the parameter. Thus $E$ is a continuous simple eigenvalue
branch, separated by at least $\Delta/2$ from every other eigenvalue.
At zero it is the ground eigenvalue of $H_0$. Continuity of the
finite-dimensional eigenvalues prevents a crossing, proving
\eqref{sc:gap}.
\end{proof}

\begin{proof}[Extension to compact-resolvent site operators]
Let $P_{u,R}=\boldsymbol{1}_{[0,R]}(h_u)$ and
$\mathcal F_R=\bigotimes_u\Ran P_{u,R}$. These are finite-dimensional
spaces containing the product vacuum. Compress $H$ to $\mathcal F_R$.
Each onsite gap remains at least $\Delta$, and each compressed
interaction is supported on the same $p$ and has norm at most $J$.
The finite-dimensional theorem therefore applies with the same constants.

In a finite system $V$ is bounded. Hence $H$ is self-adjoint on
$\Dom H_0$, bounded below, and has compact resolvent. The nested
union of the spaces $\mathcal F_R$ contains all finite linear
combinations of product eigenvectors of the $h_u$, and is an operator
core, hence also a form core, for $H_0$ and $H$.
The min--max principle gives, for each fixed eigenvalue index $j$,
\[
 E_j(H|_{\mathcal F_R})\mathbin{\downarrow}E_j(H),
\]
where $H|_{\mathcal F_R}$ denotes compression and eigenvalues are
ordered with multiplicity. If the full space has dimension at least
two, pass to the limit in
$E_1(H|_{\mathcal F_R})-E_0(H|_{\mathcal F_R})\ge\Delta/2$.
The resulting positive difference proves both the asserted gap and
simplicity. For a one-dimensional full space, the spectral exclusion
is vacuous and the ground state is simple. No norm of an unbounded
onsite operator is used in the estimates.
\end{proof}

\subsection{Explicit rational certificates}
\label{sc:certificates}

The following choices make the criterion explicit:
\begin{equation}
 \lambda_3=0.01790402,\qquad \lambda_4=0.01190829.
 \label{sc:certified-lambdas}
\end{equation}
Every terminating decimal in this subsection denotes an exact rational
number. To check the inequalities without locating a polynomial root,
use the following rational upper bounds and witnesses:
\begin{center}
\begin{tabular}{c c c c}
$s$ & $\overline\gamma_0$ & $\overline\gamma_1$ & $r$\\
\hline
3 & 1.269295517644 & 5.859465277083 & 0.03911425\\
4 & 1.464866319271 & 8.597480251020 & 0.02996979
\end{tabular}
\end{center}
The exact squares required by \eqref{sc:gammas} are
\begin{equation}
 \begin{array}{c|cc}
 s&\gamma_0(s)^2&\gamma_1(s)^2\\ \hline
 3&29/18&103/3\\
 4&103/48&887/12
 \end{array}.
 \label{sc:exact-gamma-squares}
\end{equation}
Squaring the displayed rational upper bounds verifies
$\overline\gamma_i^2\ge\gamma_i^2$.
Let $\overline q_s$ be \eqref{sc:majorant} with these upper bounds
substituted in \eqref{sc:betas}. Direct rational substitution gives,
for each of the two rows,
\begin{equation}
 r-\lambda_s\overline q_s(r)>\frac{3}{10^8},
 \qquad
 1-2\lambda_s\overline q_s'(r)>\frac{9}{10^7}.
 \label{sc:rational-slacks}
\end{equation}
These are finite rational inequalities, verifiable by clearing
denominators. Since $q_s$ and $q_s'$ have nonnegative coefficients
bounded by those of $\overline q_s$, they imply
\eqref{sc:criterion} for all $0\le\lambda\le\lambda_s$.
The certificate checks only this arithmetic; the universal operator
conclusion comes from the preceding proof. It does not follow from
finite matrix diagonalizations or from sampling interaction operators.

\subsection{Application to the specified $\SU(N)$ lattice Hamiltonian}

Let $d\ge2$ be the spatial dimension, $N\ge2$, $x=g^2>0$, and
$0\le b_N\le2N$. On a finite collection $E$ of positively oriented
cubic-lattice links, retain any subset $\mathcal P$ of the usual
elementary plaquettes whose four links belong to $E$, with no repetitions
of an anchored plaquette. Reverse
orientations denote inverses of the positive link variables.
At lattice spacing one, use
\begin{equation}
 H_{\rm lat}(x)=\frac{x}{2}\sum_{l\in E}C_l+
       \frac{b_N}{x}\sum_{p\in\mathcal P}(1-T_p),
 \qquad T_p=\frac{\operatorname{Re}\Tr U_p}{N}.
 \label{sc:KS-H}
\end{equation}
The link space is $L^2(\SU(N),\dd U)$ for normalized Haar measure;
$C_l$ is the positive group Casimir with
$\Tr(t_at_b)=\delta_{ab}/2$. The physical space is invariant under
the vertex gauge groups, with no external charges. These conventions
specify the magnetic normalization explicitly; compare the Hamiltonian
formulation in \cite{KS}. The bounds below concern the dimensionless
energies of \eqref{sc:KS-H}.

The least nonzero Casimir is
\begin{equation}
 C_F=\frac{N^2-1}{2N}\ge\frac34.
 \label{sc:CF}
\end{equation}
For completeness, write a nonzero dominant highest weight as
$\Lambda=\omega_i+\nu$, where $\nu$ is dominant and the roots have
squared length two. The Gram entries
$(\omega_i,\omega_j)=\min(i,j)(N-\max(i,j))/N$ are positive, so
$C_2(\Lambda)\ge C_2(\omega_i)$ for
$C_2(\Lambda)=(\Lambda,\Lambda+2\rho)/2$,
where $\rho=\sum_{j=1}^{N-1}\omega_j$.
Since $C_2(\omega_i)=i(N-i)(N+1)/(2N)\ge C_F$, and the
fundamental representation attains equality, Peter--Weyl decomposition
proves \eqref{sc:CF} on the entire link space.

Divide \eqref{sc:KS-H} by $s_0=xC_F/2$ and remove its scalar
magnetic constant. Precisely,
\begin{equation}
 \frac{H_{\rm lat}}{s_0}-\varepsilon|\mathcal P|I
   =\sum_{l\in E}\frac{C_l}{C_F}
        +\varepsilon\sum_{p\in\mathcal P}(-T_p),
 \qquad
 \varepsilon=\frac{2b_N}{x^2C_F}\le\frac{32}{3x^2}.
 \label{sc:normalization}
\end{equation}
The last inequality uses
$2b_N/C_F\le8N^2/(N^2-1)\le32/3$.
The normalized onsite gap is $\Delta=1$, $J=1$ because
$\|T_p\|\le1$, and each plaquette meets $s=4$ link sites.
Each link belongs to at most $D=2(d-1)$ elementary plaquettes.
These statements hold for any finite open subgraph, and also for
the usual periodic directed cubic lattices of side lengths at least
two, with positive-direction link factors distinguished.

\begin{corollary}[Explicit $N$-uniform finite-volume gauge gap]
\label{sc:gauge-corollary}
For the model \eqref{sc:KS-H}, let
\begin{equation}
 X_d^{\rm link}
  =\left(\frac{(32/3)\,2(d-1)}{0.01190829}\right)^{1/2}.
 \label{sc:X}
\end{equation}
If $x\ge X_d^{\rm link}$, the ground state is unique and gauge
invariant, and both on the full link space and on its physical
subspace,
\begin{equation}
 H_{\rm lat}-E_0\ge\frac{xC_F}{4}\ge\frac{3x}{16}
 \quad\hbox{on the orthogonal complement of the ground state}.
 \label{sc:gauge-gap}
\end{equation}
In three spatial dimensions the sufficient condition $x\ge59.9$
implies \eqref{sc:gauge-gap}, for every $N\ge2$ and
$0\le b_N\le2N$, uniformly over the stated finite volumes.
\end{corollary}
\begin{proof}
Equations \eqref{sc:normalization}--\eqref{sc:X} give
$\lambda=\varepsilon D\le\lambda_4$.
Theorem~\ref{sc:stability} and the certificate give a normalized
gap at least $1/2$. Restoring $s_0$ gives $xC_F/4$; adding back
the scalar constant changes no gap. All Casimirs and plaquette
traces commute with the finite product of vertex gauge groups.
The unique ground line therefore carries a continuous one-dimensional
representation of a product of $\SU(N)$ groups. Such a representation
is trivial, so the ground vector is gauge invariant. The physical
space reduces $H_{\rm lat}$ and contains its ground state;
restricting the excited-space inequality cannot lower the gap.
Finally, the exact rational inequality
$(59.9)^2(0.01190829)>(32/3)\,4$ verifies the stated threshold.
\end{proof}

For reference, upward-rounded sufficient thresholds using link sites
are $42.4$, $59.9$, and $73.4$ in spatial dimensions $2$, $3$, and $4$,
respectively. The exact values from \eqref{sc:X} are approximately
$42.326$, $59.858$, and $73.310$; these approximations are not used
as downward-rounded sufficient bounds.
Alternatively, group the $d$ outgoing positive links at each vertex
into one site. Each plaquette then meets three sites, and a vertex
site meets at most $3\binom d2$ plaquettes: for each coordinate pair
its three possible roles are the lower anchor and its two forward
neighbors. With $s=3$, $D=3d(d-1)/2$, and $\lambda_3$ from
\eqref{sc:certified-lambdas}, the sufficient threshold is
$\bigl((32/3)D/\lambda_3\bigr)^{1/2}$.
In $d=2$, this permits the slightly smaller upward-rounded value
$42.3$. Link grouping gives the better displayed bound in $d=3$.

\begin{remark}[Limits of the conclusion]
\label{sc:scope}
The uniformity established here is over finite volume and $N$ at
fixed $x$. The proof does not construct an infinite-volume GNS
Hamiltonian at the explicit threshold, nor does it prove a continuum
limit or a positive physical continuum mass. An infinite-volume
strong-coupling construction with a non-explicit threshold is a
separate consequence of the stability theorem in \cite{Yarotsky},
as applied in \cite{Spectral}. A limit with fixed 't~Hooft coupling
$Nx$ eventually leaves the present window as $N$ increases.
\end{remark}

\subsection{A counting correction in the cited BDL preprint}
\label{sc:BDL-correction}

This subsection concerns the exact version
arXiv:0707.1894v1 of \cite{BDL}, Appendix~A, p.~25, in the
paragraph preceding its equation~(69). The intermediate inference
$M_j\subset M\cup\{w\}$, and hence $|M_j|\le2|M|$, does not
follow from the support conditions stated there. The following
example also verifies that this is an operator issue, rather than
only a set-theoretic possibility.

Take three distinct qubit sites $u,v,w$, their vacuum $|000\rangle$,
and
\[
 A=a_u^\dagger a_v^\dagger a_w^\dagger,
 \qquad V_{vw}=X_vX_w,
 \qquad a^\dagger=|1\rangle\langle0|,
 \quad a=|0\rangle\langle1|.
\]
Here $X$ interchanges $|0\rangle$ and $|1\rangle$.
Direct application gives
\begin{equation}
 AV_{vw}|000\rangle=0,\qquad
 V_{vw}A|000\rangle=|100\rangle,\qquad
 \langle000|a_u[A,V_{vw}]|000\rangle=-1.
 \label{sc:BDL-counterexample}
\end{equation}
Thus a nonzero term has $M_j=\{u,v,w\}$ and $M=\{u\}$.
It obeys the source's support conditions but violates
$|M_j|\le2|M|$.

The valid general estimate is
\begin{equation}
 M_j\subset M\cup p,\quad |p|\le s,\quad M\ne\varnothing
 \quad\Longrightarrow\quad
 |M_j|\le|M|+s\le(s+1)|M|.
 \label{sc:corrected-count}
\end{equation}
For two-body interactions this gives $3|M|$, which is sharp by
\eqref{sc:BDL-counterexample}. It does not invalidate the conservative
constant $2^{13}$ in the source's Lemma~3. Repeating its counting
with the factor three instead of two gives, for the possible
commutator orders $1\le k\le4$, the coefficient bound
\begin{equation}
 k\cdot4\cdot3\cdot2^k2^{k-1}
        +4\cdot2^k2^k
 \le7168<8192=2^{13}.
 \label{sc:BDL-constant-survives}
\end{equation}
Higher orders vanish, and the zeroth-order estimate is smaller.
Accordingly the correction repairs this intermediate counting step
while preserving that theorem's stated conservative constant.
No priority claim is made for identifying the correction, and no
claim about revisions other than the specified source version is
needed here.

% SOURCE sections/vacuum.tex
\section{Fixed-block coercivity in the actual vacuum}
\label{vac:section}

This section proves a lower bound for the energy of fluctuations in a fixed
finite set of links, conditional on all exterior links, in the actual
finite-volume ground state.  The constant is independent of the number of
exterior links.  Its dependence on the number of selected links is explicit
and does not give a uniform bound for an extensive elimination layer.
We use a Hamiltonian of Kogut--Susskind type~\cite{KS} with the normalization
specified below.  The argument is presented in full; the associated research
record is available in~\cite{Archive}.

\subsection{Finite-volume Hamiltonian and its ground-state form}
\label{vac:setup}

Let $E$ be a finite set of oriented link representatives of a finite spatial
graph, let $\mathcal P$ be a finite collection of plaquette terms, and fix
$N\geq2$, $x>0$, $a>0$, and $b\geq0$.  Put
$M=\SU(N)^E$ and equip $M$ with normalized product Haar measure $\mathrm dU$.
For Hermitian generators satisfying
$\operatorname{Tr}(t_\alpha t_\beta)=\delta_{\alpha\beta}/2$, choose the
bi-invariant metric for which the corresponding left-invariant vector fields
$X_\alpha$ are orthonormal.  On link $\ell$, write
$C_\ell=-\sum_\alpha X_{\ell,\alpha}^2=-\Delta_\ell$.  Consider
\begin{equation}
 H=\frac{x}{2}\sum_{\ell\in E}C_\ell+V,
 \qquad
 V(U)=\frac b x\sum_{p\in\mathcal P}
       \left(1-\frac{\operatorname{Re}\operatorname{Tr}U_p}{N}\right).
 \label{vac:H}
\end{equation}
Here $U_p$ is the ordered plaquette holonomy.  Each summand in parentheses
lies in $[0,2]$.  The operator $H$ is the self-adjoint bounded-potential
perturbation of the product Laplacian, with operator domain $H^2(M)$ and
form domain $H^1(M)$.  The parameter $a$ enters the physical energy through
division by $a$; it has not been absorbed into $H$.

The vertex gauge group acts by
$U_\ell\mapsto g_{s(\ell)}U_\ell g_{t(\ell)}^{-1}$.  We assume the
plaquette terms and boundary conditions preserve this action.  Physical
functions are invariant under this group.  No gauge fixing is used.

\begin{lemma}[Actual vacuum and ground-state transform]
\label{vac:groundstate}
The lowest eigenvalue $E_0$ of $H$ is simple.  Its normalized eigenfunction
$\psi$ can be chosen smooth, strictly positive, and gauge invariant.  Set
$\mathrm d\mu=\psi^2\,\mathrm dU$.  Multiplication by $\psi$ is a unitary map
from $L^2(\mu)$ to $L^2(\mathrm dU)$, and the nonnegative self-adjoint operator
\begin{equation}
 K=\psi^{-1}(H-E_0)\psi/a
 \label{vac:K}
\end{equation}
has closed quadratic form
\begin{equation}
 k[f]=\frac{x}{2a}\int_M
          \sum_{\ell,\alpha}|X_{\ell,\alpha}f|^2\,\mathrm d\mu,
 \qquad f\in H^1(M).
 \label{vac:groundform}
\end{equation}
Its kernel is the space of constant functions.  These statements also hold
after restriction to the gauge-invariant subspaces.
\end{lemma}

\begin{proof}
The compact manifold $M$ is connected, whether or not the underlying graph
is connected.  Its heat kernel is strictly positive at every positive time.
The Feynman--Kac weight for the bounded real potential $V$ is strictly
positive, so $\exp(-tH)$ is positivity improving.  Compact resolvent gives a
lowest eigenfunction.  The absolute-value inequality for the Dirichlet form
allows a nonnegative minimizer, which is strictly positive by the strong
maximum principle.  Positivity improvement implies simplicity: a second
real ground eigenfunction orthogonal to the positive one would change sign,
contradicting the strict positivity statement.  Elliptic regularity gives
smoothness.  Gauge transformations commute with $H$ and preserve positivity
and norm; uniqueness therefore makes $\psi$ gauge invariant.

For smooth complex $f$, integration by parts and $H\psi=E_0\psi$ give
\begin{equation}
 \langle\psi f,(H-E_0)\psi f\rangle_{L^2(\mathrm dU)}
   =\frac{x}{2}\int_M
       \sum_{\ell,\alpha}|X_{\ell,\alpha}f|^2\psi^2\,\mathrm dU.
 \label{vac:transformidentity}
\end{equation}
Indeed, the terms containing $|f|^2|\nabla\psi|^2$ and the mixed derivatives
combine, after integration by parts, with $V|\psi f|^2$ to give
$E_0|\psi f|^2$.  On this fixed finite graph, $\psi$ is smooth and bounded
away from zero.  Multiplication by $\psi$ and its inverse preserve $H^1(M)$,
so the identity extends by form closure.  Unitarity follows directly from
$\mathrm d\mu=\psi^2\mathrm dU$.  The zero-energy functions have zero gradient
and are constant because $M$ is connected.  Multiplication by $\psi$
intertwines the gauge actions, which proves the final assertion.
\end{proof}

In particular, all energy estimates below subtract $E_0$, the actual
interacting ground energy.  No bound on $E_0$ uniform in volume is assumed.
The equivalence of weighted and unweighted Sobolev norms just used is needed
only for each fixed graph; its constants need not be uniform in volume.

\subsection{Conditional projection and finite-block theorem}
\label{vac:block}

Fix a nonempty subset $B\subseteq E$ with $m=|B|$, and write a configuration
as $(u,y)\in\SU(N)^B\times\SU(N)^{E\setminus B}$.  Let $p_B$ be the number
of plaquette terms that involve at least one link in $B$, counting each term
once.  Define
\begin{equation}
 Z(y)=\int\psi(v,y)^2\,\mathrm dv,
 \qquad
 \mathrm d\mu_y(u)=\frac{\psi(u,y)^2}{Z(y)}\,\mathrm du,
 \qquad
 Pf(u,y)=\int f(v,y)\,\mathrm d\mu_y(v),
 \label{vac:conditional}
\end{equation}
where $\mathrm du$ and $\mathrm dv$ are normalized product Haar measures.
Thus $P=\mathbb E_\mu[\,\cdot\mid y]$ is an orthogonal projection in
$L^2(\mu)$ and $P1=1$.  Put $Q=I-P$.

The restriction of the form $k$ to
$H^1(M)\cap\operatorname{Ran}Q$ is closed and densely defined in
$\operatorname{Ran}Q$; the density assertion is verified in the proof below.
Let $D_B$ be its associated self-adjoint operator on that subspace.  Thus
$D_B$ denotes a form compression.  It does not require the unbounded
operator product $QKQ$ to be defined on an assumed common operator domain.

Write
\begin{equation}
 C_F=\frac{N^2-1}{2N},\qquad n=N^2-1,
 \qquad
 \kappa_N(s)=
 \frac{\max_{u,v\in\SU(N)}k_s(u,v)}
      {\min_{u,v\in\SU(N)}k_s(u,v)},\quad s>0,
 \label{vac:constants}
\end{equation}
where $k_s$ is the one-link heat kernel of $\exp(-sC)=\exp(s\Delta)$
relative to normalized Haar measure.  Compactness and strict positivity make
$\kappa_N(s)$ finite.

\begin{theorem}[Exterior-volume-uniform finite-block coercivity]
\label{vac:finiteblock}
For every $s>0$,
\begin{equation}
 D_B\ \geq\ d_B(s)I,
 \qquad
 d_B(s)=\frac{xC_F}{2a}\,
       \kappa_N(s)^{-2m}
       \exp\!\left(-\frac{8sbp_B}{x^2}\right).
 \label{vac:heatbound}
\end{equation}
The bound is independent of the number of links outside $B$ when
$m,p_B,N,x,b,a$ are fixed.  The projection $P$ commutes with the gauge
action, so the same bound holds on the gauge-invariant part of
$\operatorname{Ran}Q$, if that part is nonzero.  If $bp_B=0$, the sharper
bound
\begin{equation}
 D_B\geq\frac{xC_F}{2a}I
 \label{vac:freebound}
\end{equation}
holds.
\end{theorem}

\begin{proof}
Split the potential into terms that avoid and meet $B$:
\begin{equation}
 V(u,y)=V_{\mathrm{out}}(y)+V_{\mathrm{hit}}(u,y),
 \qquad
 0\leq V_{\mathrm{hit}}\leq W_B:=\frac{2bp_B}{x}.
 \label{vac:split}
\end{equation}
Let $(U_t,Y_t)$ be independent link Brownian motions, with generator
$(x/2)\Delta$ on each link.  At any time $t>0$, the eigenfunction identity
and Feynman--Kac formula give
\begin{equation}
 e^{-tE_0}\psi(u,y)=
 \mathbb E_{u,y}\!\left[
 e^{-\int_0^t V(U_r,Y_r)\,\mathrm dr}\psi(U_t,Y_t)\right].
 \label{vac:fk}
\end{equation}
Remove only $V_{\mathrm{hit}}$ from this expectation and call the result
$F_t(u,y)$.  Positivity gives
\begin{equation}
 e^{-tW_B}F_t(u,y)\leq e^{-tE_0}\psi(u,y)\leq F_t(u,y).
 \label{vac:fksandwich}
\end{equation}
Put $s=xt/2$.  Independence of the link motions and the absence of $u$ from
$V_{\mathrm{out}}$ imply
\begin{align}
 F_t(u,y)&=\int k_s^B(u,v)A_t(y,v)\,\mathrm dv,
 \label{vac:factorization}\\
 A_t(y,v)&=\mathbb E_y\!\left[
       e^{-\int_0^t V_{\mathrm{out}}(Y_r)\,\mathrm dr}
       \psi(v,Y_t)\right]>0,
 \notag
\end{align}
where $k_s^B$ is the product of the $m$ one-link kernels.  For a fixed $y$,
the positive factor $A_t(y,v)$ is the same for every initial block
configuration $u$.  Consequently, for every $u,u',y$,
\begin{equation}
 \frac{\psi(u,y)}{\psi(u',y)}
 \leq e^{tW_B}\frac{F_t(u,y)}{F_t(u',y)}
 \leq e^{tW_B}\kappa_N(s)^m=:R_B(s).
 \label{vac:oscillation}
\end{equation}
The factor $e^{-tE_0}$ cancels in this same-time ratio.  The exterior
potential remains inside the common positive expectation in
\eqref{vac:factorization}; it has not been replaced by a bound depending on
the number of exterior plaquettes.  This is the source of the asserted
exterior-volume independence.

We next recall the product Haar Poincar\'e inequality in the chosen metric:
\begin{equation}
 \operatorname{Var}_{\mathrm du}(f)
 \leq\frac1{C_F}\int\sum_{\ell\in B,\alpha}
                 |X_{\ell,\alpha}f|^2\,\mathrm du.
 \label{vac:haar}
\end{equation}
For completeness, Peter--Weyl decomposition reduces this statement to the
smallest nonzero Casimir eigenvalue of $\SU(N)$.  With long roots of squared
length $2$, the Casimir of highest weight $\lambda$ is
$(\lambda,\lambda+2\rho)/2$.  The fundamental weight $\omega_j$ has Casimir
$j(N-j)(N+1)/(2N)$.  All pairings of fundamental weights are nonnegative,
so every nonzero dominant integral weight has Casimir at least the minimum
of these fundamental values, namely $C_F$.  On a product, the eigenvalues
add; a nonconstant component therefore has eigenvalue at least $C_F$.
This proves~\eqref{vac:haar}, including for complex functions by applying
the real inequality to their real and imaginary parts.

For any positive probability density $w$ relative to product Haar with
$\max w/\min w\leq R^2$, the minimizing-constant definition of variance and
\eqref{vac:haar} give
\begin{align}
 \operatorname{Var}_{w\mathrm du}(f)
 &\leq (\max w)\operatorname{Var}_{\mathrm du}(f)\notag\\
 &\leq\frac{\max w}{C_F}\int|\nabla_Bf|^2\,\mathrm du
 \leq\frac{R^2}{C_F}\int|\nabla_Bf|^2w\,\mathrm du,
 \label{vac:densitycomparison}
\end{align}
where $|\nabla_Bf|^2=\sum_{\ell\in B,\alpha}|X_{\ell,\alpha}f|^2$.
No Haar-mean-zero condition is imposed on $f$ in this comparison.
By~\eqref{vac:oscillation}, the density of $\mu_y$ in
\eqref{vac:conditional} has maximum-to-minimum ratio at most $R_B(s)^2$.
Applying~\eqref{vac:densitycomparison} conditionally and integrating over
the exterior marginal yields
\begin{equation}
 \|Qf\|_{L^2(\mu)}^2
 \leq\frac{R_B(s)^2}{C_F}
        \int_M|\nabla_Bf|^2\,\mathrm d\mu.
 \label{vac:conditionalpoincare}
\end{equation}
If $f\in\operatorname{Ran}Q$, the exact ground-state form therefore gives
\begin{equation}
 k[f]\geq\frac{x}{2a}\int_M|\nabla_Bf|^2\,\mathrm d\mu
       \geq\frac{xC_F}{2aR_B(s)^2}\|f\|_{L^2(\mu)}^2.
 \label{vac:discardedform}
\end{equation}
Since $t=2s/x$ and $W_B=2bp_B/x$, this is precisely
\eqref{vac:heatbound}.

To justify the form restriction, observe that $P$ maps smooth functions to
smooth functions.  On each fixed compact graph it is bounded on $H^1(M)$:
differentiating its conditional integral in an exterior direction produces
the conditional derivative of $f$ and a covariance with
$2\nabla_y\log\psi$.  These coefficients are smooth and bounded, and
conditional Cauchy--Schwarz gives the required fixed-graph bound.  Derivatives
in the $B$ directions vanish for $Pf$.  Thus $Q$ is also bounded on
$H^1(M)$.  If $f\in H^1(M)\cap\operatorname{Ran}Q$, approximate $f$ in
$H^1$ by smooth $f_j$ and use $Qf_j\to f$.  This proves smooth density in
the restricted form domain.  Smooth density in $L^2$ and boundedness of $Q$
also prove that the restriction is densely defined in
$\operatorname{Ran}Q$.  It is closed because this subspace is closed in
$L^2(\mu)$ and $k$ is closed.  Inequality~\eqref{vac:discardedform} consequently
extends to the entire restricted form domain without any uniform derivative
bound on $\psi$.

For a fixed gauge transformation, the selected and exterior link coordinates
transform separately.  Both transformations preserve their Haar measures,
and $\psi$ is invariant.  Changing variables in
\eqref{vac:conditional} therefore shows that $P$ intertwines the gauge
action.  The forms and their restrictions are gauge invariant, proving the
physical-sector assertion.

Finally, if $bp_B=0$, the Hamiltonian splits as a free product Laplacian on
$B$ plus an exterior operator.  Its positive ground function is constant on
$B$.  The conditional measures are exactly product Haar, and
\eqref{vac:haar} gives~\eqref{vac:freebound}.
\end{proof}

\subsection{An explicit estimate at weak coupling}
\label{vac:explicit}

Let $D_N$ be the diameter of $\SU(N)$ in the specified Casimir metric.
The next heat-kernel estimate supplies explicit constants in
Theorem~\ref{vac:finiteblock}.

\begin{lemma}[Compact heat-kernel ratio]
\label{vac:heatlemma}
For every $s>0$,
\begin{equation}
 \kappa_N(s)\leq 2^n
       \exp\!\left(\frac{3D_N^2}{4s}\right),
 \qquad n=N^2-1.
 \label{vac:kappabound}
\end{equation}
\end{lemma}

\begin{proof}
The bi-invariant metric has nonnegative Ricci curvature.  For a positive
solution of $u_t=\Delta u$ on a compact $n$-dimensional manifold with
nonnegative Ricci curvature, the Li--Yau estimate~\cite{LiYau} is
\begin{equation}
 |\nabla\log u|^2-\partial_t\log u\leq\frac{n}{2t}.
 \label{vac:liyau}
\end{equation}
Here is its compact, boundaryless derivation.  Set $f=\log u$,
$w=|\nabla f|^2-f_t$, and $\mathcal F=tw$.  The heat equation, Bochner
identity, and $|\operatorname{Hess}f|^2\geq(\Delta f)^2/n$ imply
\begin{equation}
 (\Delta-\partial_t)\mathcal F
 \geq\frac{2\mathcal F^2}{nt}-\frac{\mathcal F}{t}
       -2\langle\nabla f,\nabla\mathcal F\rangle.
 \label{vac:bochnerheat}
\end{equation}
For smooth positive initial data, $\mathcal F=0$ at $t=0$.  A positive
space-time maximum exceeding $n/2$ contradicts~\eqref{vac:bochnerheat}
and the maximum principle.  This proves~\eqref{vac:liyau}.  For the heat
kernel, start after a positive time shift and then let that shift tend to
zero.  Compactness removes boundary and escape issues.

Along a constant-speed minimizing geodesic $\gamma$ between times
$0<t_1<t_2$, inequality~\eqref{vac:liyau} gives
\[
 \frac{\mathrm d}{\mathrm dt}f(\gamma(t),t)
 \geq-\frac{|\dot\gamma(t)|^2}{4}-\frac{n}{2t}.
\]
Integration yields
\begin{equation}
 u(z,t_1)\leq u(w,t_2)
       \left(\frac{t_2}{t_1}\right)^{n/2}
       \exp\!\left(\frac{\operatorname{dist}(z,w)^2}
                         {4(t_2-t_1)}\right).
 \label{vac:harnack}
\end{equation}
Apply this to $k_t(v,\cdot)$ between $s$ and $2s$ and integrate the later
point $w$ against normalized Haar measure.  Total heat-kernel mass one gives
\begin{equation}
 \max k_s\leq 2^{n/2}e^{D_N^2/(4s)}.
 \label{vac:heatupper}
\end{equation}
Apply~\eqref{vac:harnack} between $s/2$ and $s$ and integrate the earlier
point $z$ instead.  The result is
\begin{equation}
 \min k_s\geq2^{-n/2}e^{-D_N^2/(2s)}.
 \label{vac:heatlower}
\end{equation}
The two inequalities prove~\eqref{vac:kappabound}.  Normalizing Riemannian
volume to Haar probability multiplies the heat kernel by a constant and
does not change the logarithmic differential estimate; its mass in the
normalized measure is one, as used above.
\end{proof}

\begin{corollary}[Explicit finite-block lower bound]
\label{vac:exponential}
If $bp_B>0$, then
\begin{equation}
 D_B\geq\frac{xC_F}{2a}\,2^{-2mn}
                   e^{-C_B/x}I,
 \qquad C_B=4D_N\sqrt{3mbp_B}.
 \label{vac:explicitbound}
\end{equation}
\end{corollary}

\begin{proof}
Substituting~\eqref{vac:kappabound} into~\eqref{vac:heatbound} gives the
lower certificate
\[
 \frac{xC_F}{2a}\,2^{-2mn}
 \exp\!\left(-\frac{3mD_N^2}{2s}-\frac{8sbp_B}{x^2}\right).
\]
The positive exponent is minimized at
\[
 s_* =\frac{xD_N}{4}\sqrt{\frac{3m}{bp_B}},
 \qquad
 \frac{3mD_N^2}{2s_*}+\frac{8s_*bp_B}{x^2}
       =\frac{4D_N\sqrt{3mbp_B}}{x}.
\]
This proves~\eqref{vac:explicitbound}.  The case $bp_B=0$ is already covered
by~\eqref{vac:freebound} and does not require this optimization.
\end{proof}

The exponential dependence is a rigorous lower certificate, not an
asymptotic formula for the actual discarded spectrum as $x\downarrow0$.
In particular, the estimate is not uniform in $m$ or $N$.

\subsection{Conditional scale budget and the remaining limitation}
\label{vac:budget}

One can quantify the strength of this certificate without asserting a
renormalization theorem.  Suppose, as an additional hypothesis, that the same
fixed-block estimate with constants
$c_0=(C_F/2)2^{-2mn}$ and $C=C_B$ were available for every relevant exact
coarse operator.  Consider the illustrative scale and coupling sequences
\begin{equation}
 a_r=A L^{-r},\qquad x_r=(\sigma+qr)^{-1},\qquad r\geq1,
 \quad A>0,\ L>1,\ \sigma>0,\ q\geq0.
 \label{vac:modelscales}
\end{equation}
The reciprocal costs of the certified bounds
$d_r^{\mathrm{cert}}=c_0x_r e^{-C/x_r}/a_r$ then satisfy the exact identity
\begin{equation}
 \sum_{r\geq1}\frac1{d_r^{\mathrm{cert}}}
 =\frac{A}{c_0}e^{C\sigma}
       \sum_{r\geq1}(\sigma+qr)\rho^r,
 \qquad \rho=\frac{e^{Cq}}{L}.
 \label{vac:budgetseries}
\end{equation}
This series is finite exactly when $\rho<1$, in which case it equals
\begin{equation}
 \frac{A}{c_0}e^{C\sigma}
 \left(\frac{\sigma\rho}{1-\rho}
       +\frac{q\rho}{(1-\rho)^2}\right).
 \label{vac:budgetvalue}
\end{equation}
No actual Yang--Mills coupling trajectory, choice of scale factor, or
stability of the finite-block theorem under exact elimination is proved by
this calculation.  If $\rho\geq1$, this particular certificate does not
give a finite upper bound on the inverse-gap cost; the true discarded gaps
could still give a finite sum.

Theorem~\ref{vac:finiteblock} concerns one fixed finite block in the original
Hamiltonian~\eqref{vac:H}.  A whole elimination layer generally contains an
extensive number of links.  If that layer is used as $B$, the factor
$2^{-2mn}$ deteriorates exponentially in $m$, and $C_B$ grows proportionally
to $m$ when $p_B$ does.  Eliminating bounded blocks one after another instead
requires estimates for the successive exact coarse operators and control of
the number of elimination steps.  Neither follows from the theorem: the
generated operators have not been shown to retain the form~\eqref{vac:H},
and conditional bounds for separate blocks do not by themselves provide a
collective conditional Poincar\'e inequality in an interacting vacuum.
Thus the proved exterior-volume uniformity leaves open a whole-layer
estimate for the actual compact-link ground measure that persists under
exact coarse graining.
